% =============================================================================
%  Grupos y Anillos --- Notas del Examen II
%  Compilar con: pdflatex (dos veces, para que el indice recoja cada resultado)
%  Plantilla basada en las notas de Analisis I del mismo autor.
% =============================================================================
\documentclass[11pt,twoside]{article}

\usepackage[margin=0.9in]{geometry}
\usepackage[utf8]{inputenc}
\usepackage[T1]{fontenc}
% --- Fuentes: usa XCharter+newtxmath si estan instaladas; si no, degrada ---
\IfFileExists{XCharter.sty}{%
  \usepackage{XCharter}%
  \IfFileExists{newtxmath.sty}{\usepackage[libertine,charter]{newtxmath}}{}%
}{%
  \IfFileExists{charter.sty}{\usepackage{charter}}{\usepackage{lmodern}}%
}
% --- Idioma: usa babel-spanish si esta disponible; si no, ingles o nada ---
\IfFileExists{babel.sty}{%
  \IfFileExists{spanish.ldf}{\usepackage[spanish,es-noshorthands]{babel}}{%
    \IfFileExists{english.ldf}{\usepackage[english]{babel}}{}}%
}{}
\usepackage{amsmath,amssymb,amsthm,mathtools}
\usepackage{enumitem}
\usepackage[protrusion=true,expansion=false]{microtype}
\usepackage{array}
\usepackage{longtable}
\usepackage{titlesec}
\usepackage{xcolor}
\usepackage[most]{tcolorbox}
\usepackage{tikz}
\usetikzlibrary{positioning, arrows.meta, shapes.geometric, calc, fit}
\usepackage{eso-pic}
\usepackage{etoolbox}
\usepackage{xstring}
\usepackage{fancyhdr}
\usepackage{lastpage}
\usepackage{hyperref}

\hypersetup{
  colorlinks=true,
  linkcolor=blue!50!black,
  urlcolor=blue!50!black,
  pdftitle={Grupos y Anillos --- Notas del Examen II},
  pdfauthor={Sebastián Fernández Rivera}
}

\AtBeginDocument{%
  \renewcommand{\contentsname}{\'Indice}%
  \renewcommand{\refname}{Referencias}%
  \renewcommand{\proofname}{\textbf{Prueba}}%
}

\setcounter{tocdepth}{3}
\setcounter{secnumdepth}{2}
\allowdisplaybreaks

\newcounter{result}[subsection]
\renewcommand{\theresult}{\thesubsection.\arabic{result}}

\newenvironment{namedres}[2]{%
  \refstepcounter{result}\phantomsection%
  \addcontentsline{toc}{subsubsection}{\protect\numberline{\theresult}#1 (\textit{#2})}%
  \par\medskip\noindent\textbf{#1 \theresult\ (\textit{#2}).}%
  \par\nopagebreak\noindent\itshape\ignorespaces%
}{\par\medskip\upshape}

\newenvironment{namedstd}[2]{%
  \refstepcounter{result}\phantomsection%
  \addcontentsline{toc}{subsubsection}{\protect\numberline{\theresult}#1 (\textit{#2})}%
  \par\medskip\noindent\textbf{#1 \theresult\ (\textit{#2}).}%
  \par\nopagebreak\noindent\ignorespaces%
}{\par\medskip}

\newenvironment{namedresx}[1]{%
  \refstepcounter{result}\phantomsection%
  \addcontentsline{toc}{subsubsection}{\protect\numberline{\theresult}#1}%
  \par\medskip\noindent\textbf{#1 \theresult.}%
  \par\nopagebreak\noindent\itshape\ignorespaces%
}{\par\medskip\upshape}

\newenvironment{namedstdx}[1]{%
  \refstepcounter{result}\phantomsection%
  \addcontentsline{toc}{subsubsection}{\protect\numberline{\theresult}#1}%
  \par\medskip\noindent\textbf{#1 \theresult.}%
  \par\nopagebreak\noindent\ignorespaces%
}{\par\medskip}


% -----------------------------------------------------------------------------
%  Cajas del Anexo (estilo de las notas de Análisis I)
% -----------------------------------------------------------------------------
\newtcolorbox{equivbox}[2]{%
  enhanced,
  breakable,
  colback=#1!6!white,
  colframe=#1!72!black,
  coltitle=white,
  fonttitle=\bfseries,
  fontupper=\small,
  title={#2},% las llaves protegen comas en el título del parser de pgfkeys
  arc=4pt,
  boxrule=0.8pt,
  left=10pt, right=10pt, top=6pt, bottom=6pt,
  before skip=4pt, after skip=8pt,
}

\newtcolorbox{strategybox}[1]{%
  enhanced,
  breakable,
  colback=gray!4,
  colframe=blue!50!black,
  coltitle=white,
  fonttitle=\bfseries,
  fontupper=\small,
  title={#1},
  arc=3pt,
  boxrule=0.6pt,
  left=8pt, right=8pt, top=4pt, bottom=4pt,
  before skip=8pt, after skip=8pt,
}

% Encabezado compacto para los sub-grupos de un Anexo
\newcommand{\anexogroup}[1]{%
  \par\addvspace{0.7\baselineskip}%
  \noindent{\bfseries\large #1}\par\nobreak
  \addcontentsline{toc}{subsection}{#1}%
  \vspace{4pt}\nopagebreak
}

\newcommand{\hlu}[1]{\underline{#1}}

\DeclareMathOperator{\Ker}{Ker}
\DeclareMathOperator{\Imm}{Im}
\DeclareMathOperator{\Aut}{Aut}
\DeclareMathOperator{\sgn}{sgn}
\DeclareMathOperator{\mcd}{mcd}
\DeclareMathOperator{\MCD}{MCD}
\newcommand{\Z}{\mathbb{Z}}
\newcommand{\R}{\mathbb{R}}
\newcommand{\Q}{\mathbb{Q}}
\newcommand{\N}{\mathbb{N}}
\newcommand{\F}{\mathbb{F}}

\pagestyle{fancy}
\fancyhf{}
\renewcommand{\sectionmark}[1]{\markboth{\thesection.\ #1}{}}
\fancyhead[L]{\small\nouppercase{\leftmark}}
\fancyhead[R]{\small\textit{Grupos y Anillos --- Examen II}}
\fancyfoot[C]{\small Página \thepage{} de \pageref{LastPage}}
\renewcommand{\headrulewidth}{0.4pt}
\renewcommand{\footrulewidth}{0pt}
\setlength{\headheight}{14pt}

\definecolor{sec1col}{HTML}{1F77B4}
\definecolor{sec2col}{HTML}{2CA02C}
\definecolor{sec3col}{HTML}{D62728}
\definecolor{sec4col}{HTML}{9467BD}
\definecolor{sec5col}{HTML}{FF7F0E}
\definecolor{sec6col}{HTML}{8C564B}
\definecolor{sec7col}{HTML}{17BECF}
\definecolor{sec8col}{HTML}{BCBD22}
\definecolor{sec9col}{HTML}{E377C2}

\newcounter{thumbsec}
\pretocmd{\section}{\stepcounter{thumbsec}}{}{}
\newif\ifdrawthumb
\drawthumbfalse

\newcommand{\drawthumbtab}{%
  \ifdrawthumb
  \ifnum\value{thumbsec}>0
    \ifnum\value{thumbsec}<10
      \begin{tikzpicture}[remember picture, overlay]%
        \pgfmathsetmacro{\toppadcm}{4.0}%
        \pgfmathsetmacro{\tabHcm}{1.6}%
        \pgfmathsetmacro{\ypos}{-\toppadcm - (\value{thumbsec} - 1) * \tabHcm}%
        \ifodd\value{page}%
          \fill[sec\arabic{thumbsec}col]
            ([xshift=0pt, yshift=\ypos cm]current page.north east)
            rectangle ++(-0.7cm, -1.2cm);
          \node[white, font=\bfseries\large, anchor=center]
            at ([xshift=-0.35cm, yshift=\ypos cm - 0.6cm]current page.north east)
            {\arabic{thumbsec}};
        \else%
          \fill[sec\arabic{thumbsec}col]
            ([xshift=0pt, yshift=\ypos cm]current page.north west)
            rectangle ++(0.7cm, -1.2cm);
          \node[white, font=\bfseries\large, anchor=center]
            at ([xshift=0.35cm, yshift=\ypos cm - 0.6cm]current page.north west)
            {\arabic{thumbsec}};
        \fi%
      \end{tikzpicture}%
    \fi
  \fi
  \fi
}
\AddToShipoutPictureFG{\drawthumbtab}

\begin{document}

\begin{titlepage}
\thispagestyle{empty}
\centering
\vspace*{1.5cm}
{\scshape\Large Universidad de Costa Rica\par}
\vspace{0.3cm}
{\scshape\large Escuela de Matem\'atica\par}
\vspace{1.2cm}
{\color{blue!50!black}\rule{\textwidth}{1.2pt}}
\vspace{0.6cm}
{\Huge\bfseries Grupos y Anillos\par}
\vspace{0.5cm}
{\Large\itshape Notas del Examen II\par}
\vspace{0.4cm}
{\color{blue!50!black}\rule{\textwidth}{1.2pt}}
\vspace{1.4cm}
{\Large\bfseries Teor\'ia de Grupos\par}
\vspace{2.2cm}
\begin{tabular}{rl}
\textbf{Autor:}                & Sebasti\'{a}n Fern\'{a}ndez Rivera \\[4pt]
\textbf{Curso:}                & Grupos y Anillos \\[4pt]
\textbf{Documento:}            & Notas para el II Examen Parcial \\[4pt]
\textbf{Resultados numerados:} & 170$+$ \\[4pt]
\textbf{Fecha:}                & \today
\end{tabular}
\vfill
\begin{minipage}{0.85\textwidth}
\centering\itshape
Este documento re\'{u}ne las definiciones, teoremas, proposiciones, lemas, corolarios,
ejemplos, notas, ejercicios y notaciones del curso. Cada resultado est\'{a} numerado de
la forma \texttt{seccion.subseccion.resultado} y aparece en el \'{\i}ndice con su
p\'{a}gina para una consulta r\'{a}pida durante el examen.
\end{minipage}
\vspace{1cm}
{\color{blue!50!black}\rule{0.5\textwidth}{0.6pt}}
\end{titlepage}

\thispagestyle{empty}
{\hypersetup{linkcolor=black}\tableofcontents}
\newpage
\setcounter{thumbsec}{0}
\drawthumbtrue


% ======================================================================
\section{Grupos}
% ======================================================================

\subsection{Definiciones elementales}

\begin{namedstd}{Definición}{Operación binaria}
Sea $G$ un conjunto no vacío. Una operación binaria es una aplicación $G \times G \to G$ con $(g_{1},g_{2}) \mapsto g_{3}$.

\end{namedstd}
\begin{namedstd}{Definición}{Grupo}
Un grupo $G$ es un conjunto no vacío junto con una operación binaria $\ast: G \times G \to G$ tal que:
\begin{enumerate}[label=(\arabic*)]
\item $\ast$ es asociativa:
\[
\forall g_{1},g_{2},g_{3} \in G \quad  (g_{1} \ast g_{2}) \ast g_{3} = g_{1} \ast (g_{2} \ast g_{3});
\]
\item existe un elemento neutro $e \in G$ tal que $e \ast g = g\ast e = g$ para todo $g \in G$. Muchas veces, escribimos $1_{G}$ para denotar el elemento neutro del grupo $G$;
\item existe un inverso para todo elemento, esto es,
\[
\forall g \in G  \quad \exists h \in G  \quad (g \ast h = h \ast g = e).
\]
\end{enumerate}

\end{namedstd}
\begin{namedstd}{Definición}{Grupo Abeliano}
Si $(G, \ast)$ es un grupo que además cumple que para todos $g_{1},g_{2} \in G$, $g_{1} \ast g_{2} = g_{2} \ast g_{1}$, decimos que $G$ es \hlu{conmutativo} o \hlu{abeliano}.

\end{namedstd}
\begin{namedstd}{Definición}{Monoide}
Un monoide es una estructura más general, dotada por un conjunto, una operación binaria asociativa con elemento neutro, pero en donde no necesariamente todo elemento tiene un inverso.
\end{namedstd}

\subsection{Propiedades de los grupos}

\begin{namedres}{Proposición}{Unicidad del neutro}
Si $e, e'$ son dos elementos neutros, $e' = e' \ast e = e$.

\end{namedres}
\begin{namedres}{Proposición}{Unicidad del inverso}
Sea $g \in G$ y sean $h, h'$ inversos de $g$. Entonces
\[
h' = e \ast h' = (h \ast g) \ast h' = h \ast (g \ast h') = h \ast e = h.
\]
\end{namedres}
\begin{namedstd}{Notación}{Inverso}
Si $(G, \ast)$ es un grupo y $g \in G$, denotamos por $g^{-1}$ al inverso de $g$.

\end{namedstd}
\begin{namedstd}{Nota}{Propiedades del inverso}
\begin{enumerate}[label=(\arabic*)]
\item $(g^{-1})^{-1} = g$ para todo $g \in G$
\item $(g \ast h)^{-1} = h ^{-1} \ast g ^{-1}$ para todos $g,h \in G$.
\end{enumerate}

\end{namedstd}
\begin{namedres}{Proposición}{Leyes de cancelación}
Si $g, h_{1}, h_{2} \in G$, entonces
\begin{enumerate}[label=(\arabic*)]
\item $g \ast h_{1} = g \ast h_{2} \implies h_{1}=h_{2}$,
\item $h_{1} \ast g = h_{2} \ast g \implies h_{1} = h_{2}$.
\end{enumerate}

\end{namedres}

\subsection{Ejemplos de grupos}

Algunos ejemplos son muy importantes dentro de la estructura

\begin{namedstd}{Ejemplo}{Grupo trivial}
Sea $G = \{  e \}$ tal que $e \ast e = e$. $(G, \ast)$ es un grupo y es llamado grupo trivial.

\end{namedstd}
\begin{namedstd}{Ejemplo}{Grupos aditivos}
$(\mathbb{Z}, + )$ es un grupo con elemento neutro $0$ y en donde el inverso de $z \in Z$ es $-z$. $(\mathbb{N}, +)$ no es grupo pero sí es monoide, pero $(\mathbb{N}^{\ast},+)$ no es monoide. $(\mathbb{R},+), (\mathbb{Q},+), (\mathbb{C}, +), (\mathbb{Z},+)$ son todos grupos abelianos.

\end{namedstd}
\begin{namedstd}{Ejemplo}{Grupos multiplicativos}
$(\mathbb{R}^{\ast}, \cdot), (\mathbb{Q}^{\ast}, \cdot), (\mathbb{C}^{\ast}, \cdot)$ son grupos abelianos con neutro $1$ y donde el inverso de $x$ es $\frac{1}{x}$. En general, si $\mathbb{F}$ es un cuerpo, $(\mathbb{F}^{\ast}, \cdot)$ es un grupo. $(\mathbb{Z}^{\ast}, \cdot)$ no es un grupo porque no todo elemento posee un inverso.

\end{namedstd}
\begin{namedstd}{Ejemplo}{Grupos lineales}
Sea $n \in \mathbb{N}^{\ast}$ y $\mathbb{F}$ un cuerpo. Definimos
\[
GL_{n}(\mathbb{F}) = \{\text{matrices }n \times n \text{ invertibles con entradas en }\mathbb{F} \}.
\]
$(GL_{n}(\mathbb{F}), \cdot)$, donde $\cdot$ es la multiplicación estándar de matrices. A este grupo se le llama el grupo lineal sobre $\mathbb{F}$.

\end{namedstd}
\begin{namedstd}{Ejemplo}{Grupos sobre congruencias}
Sea $m>0$ y considere $\mathbb{Z}_{m} = \mathbb{Z}/{\equiv_{m}}$ ($\mathbb{Z}$ cocientado en la relación de congruencia módulo $m$), i.e., $\mathbb{Z}_{m} = \{ [0], [1],\dots, [m-1] \}$. Defina la operación $\hat{+}$ en $\mathbb{Z}_{m}$ tal que $[a]\hat{+}[b] = [a+b]$ (hay que mostrar que está bien definida). $(\mathbb{Z}_{m}, \hat{+})$ es un grupo con elemento neutro $[0]$. Si $[n] \in \mathbb{Z}_{m}$, el inverso es $[-n]$. Consideremos el ejemplo con $m=2$.


Para $\mathbb{Z}_{2}$, tenemos $\mathbb{Z}_{2} = \{ [0]. [1] \}$. Observe la tabla de operaciones

\begin{center}
\begin{tabular}{|c|c|c|}\hline
$\hat{+}$ & $[0]$ & $[1]$ \\ \hline
$[0]$ & $[0]$ & $[1]$ \\ \hline
$[1]$ & $[1]$ & $[0]$ \\ \hline
\end{tabular}
\end{center}

Definamos ahora sobre $\mathbb{Z}_{m}^{\ast} = \mathbb{Z}_{m} \setminus \{ [0] \}$. Definamos la operación $\hat{\cdot}$ como
\[
[a] \hat{\cdot} [b] = [a \cdot b]
\]
Queremos construir un grupo multiplicativo con esta nueva operación. ¿Para cuales casos de $m$ está esto bien definido y existen inversos? Para $m$ compuesto, la multiplicación no está bien definida en $\mathbb{Z}_{m}^{\ast}$ porque puedo tomar elementos de ciertas clases de equivalencia cuyo producto tiene resultado $[0]$.

Para primos, la operación está bien definidos. Sea $n$ primo y asuma que existen $[a], [b] \in \mathbb{Z}_{n}^{\ast}$ tales que $[a] \hat{\cdot} [b] = [a \cdot b] = [0]$. Esto implica que $n \mid a \cdot b$, pero como $n$ es primo, $n \mid a$ o $n \mid b$, de modo que $[a] = [0]$ o $[b]=[0]$, una contradicción pues $[0] \not\in \mathbb{Z}_{n}^{\ast}$.

Probaremos que $(\mathbb{Z}_{n}^{\ast}, \hat{\cdot})$ es un grupo para $n$ primo.  Note que $\mathbb{Z}_{n}^{\ast} = \{ [1],\dots,[n-1] \}$. Sea $r \in \{ 1,\dots,n-1 \}$. Como $n$ es primo, tenemos que $\MCD(r,n) = 1$. Luego, existen enteros $q,s$ tales que $qr+sn = 1$. Esto implica que
\[
[q \cdot r] = [1] \implies [q] \hat{\cdot} [r] = [1].
\]
Podemos suponer que $q \in \{ 1,\dots,n-1 \}$, pues de lo contrario tomamos el representante de su clase de equivalencia. Así, $[q]$ es el inverso de $[r]$. Así, $(\mathbb{Z}_{n}^{\ast}. \hat{\cdot})$ tiene inversos para todo elemento, y por tanto es un grupo (las otras propiedades son triviales).

\end{namedstd}
\begin{namedstd}{Nota}{Cuerpo $\Z_p$}
Es fácil verificar que si $p$ es primo $(\mathbb{Z}_{p}, \hat{+}, \hat{\cdot}, [0], [1])$ es un cuerpo.

\end{namedstd}
\begin{namedstd}{Ejemplo}{Grupos de simetrías de un cuadrado}
Considere un cuadrado con sus vértices numerados. ¿De cuántas maneras puedo mover sus vértices para que el cuadrado siga manteniendo el mismo espacio? Puede hacer cuatro rotaciones ($0^{\circ}, 90^{\circ}, 180^{\circ}, 270^{\circ}$) en el sentido contrario a las agujas del reloj, ($r_{0},r_{90},r_{180},r_{270}$). Además, puedo hacer cuatro reflexiones: vertical ($s_{1}$), horizontal ($s_{2}$), diagonal comenzando abajo-izquierda ($s_{3}$) y diagonal comenzando arriba-izquierda ($s_{4}$).

Llamamos $D_{4}= \{ r_{0},r_{90},r_{180},r_{270}, s_{1},s_{2},s_{3},s_{4} \}$. Definimos una operación $\circ$ en $D_{4}$, como la "composición de movimientos", p.e., $r_{90} \circ s_{2}$.
\end{namedstd}
\begin{namedstd}{Ejercicio}{Grupo diédrico $D_4$}
$(D_{4},\circ)$ es un grupo pero \textbf{no} es conmutativo.

En general $D_{n}$ es el grupo de simetrías de un polígono regular de $n$ lados. Está compuesto por $r_{0},\dots,r_{n-1}$ rotaciones, donde $r_{k}$ rota la figura en un ángulo de $\frac{2\pi k}{n}$ y $s_{0}, \dots, s_{n-1}$ son $n$ reflexiones sobre las líneas $\theta = \frac{\pi k}{n}$. $D_{n}$ tiene $2n$ elemmentos, y se le llama el grupo diédrico o grupo diedral.

\end{namedstd}
\begin{namedstd}{Ejemplo}{Grupo de permutaciones de $n$ elementos}
Una permutación de $n$ elementos es una función biyectiva
\[
\begin{aligned}
\sigma:\{ 1,\dots,n \} &\to \{ 1,\dots,n \}, \text{ tal que} \\
1 &\mapsto \sigma(1)\\
& \quad\vdots \\
n &\mapsto \sigma(n).
\end{aligned}
\]
Definimos $S_{n} = \{ \sigma: \sigma \text{ es una permutación de }n \text{ elementos}\}$. $S_{n}$ es un grupo bajo la composición de funciones (ejercicio).

\end{namedstd}
\begin{namedstd}{Notación}{Permutaciones}
$\sigma: \begin{pmatrix}1 & 2 & 3 \\ 2 & 3 & 1\end{pmatrix}$ representa $\sigma:\{ 1,2,3 \} \to \{ 1,2,3 \}$, con $\sigma(1)=2, \sigma(2)=3,\sigma(3)=1$. Su inversa es $\sigma ^{-1}:\begin{pmatrix}1 & 2 & 3 \\ 3 & 1 & 2\end{pmatrix}$.

\end{namedstd}
\begin{namedstd}{Ejemplo}{Producto de grupos}
Sean $(G_{1}, \ast_{1})$ y $(G_{2}, \ast_{2})$ dos grupos. Considere el producto cartesiano $G_{1} \times G_{2}$. Defina la operación $\ast$ en $G_{1} \times G_{2}$ tal que
\[
(g_{1}, g_{2}) \ast (h_{1},h_{2}) := (g_{1} \ast_{1} h_{1}, g_{2} \ast_{2} h_{2}) \in G_{1} \times G_{2}.
\]
Entonces, $(G_{1} \times G_{2}, \ast)$ es un grupo con elemento neutro $(1_{G_{1}}, 1_{G_{2}})$. Si $(g_{1},g_{2}) \in G_{1} \times G_{2}$, entonces $(g_{1}, g_{2})^{-1} = (g_{1}^{-1}, g_{2}^{-1})$. Se puede generalizar para el producto de $n$ grupos.

\end{namedstd}

\subsection{Operaciones iteradas en grupos}

Si $(G, \ast)$ es un grupo denotamos por $g^{n} = \underbrace{ g \ast g \ast \dots \ast g }_{ n \text{ veces} }$,, $g^{0} = 1_{G}$, $g^{-n}:= (g^{-1})^{n}$.
\begin{namedres}{Proposición}{Propiedades de la potenciación}
Sea $(G, \ast)$ un grupo, $g,h \in G$. Entonces
\begin{enumerate}[label=(\arabic*)]
\item $g^{-n} = (g^{n}) ^{-1}$
\item $g^{n+m} = g^{n} \ast g^{m}$
\item Si $g \ast h = h \ast g$, entonces $(g \ast h)^{m} = g^{m} \ast h^{m}$.
\end{enumerate}

\end{namedres}

\subsection{Subgrupos}

Sea $(G, \ast)$ un grupo y $H \subseteq G$. Decimos que $H$ es un subgrupo de $G$ si $(H, \ast \mid_{H})$ es un grupo. En ese caso escribirnos $H \leq G$

\begin{namedstd}{Nota}{Verificación de subgrupo}
Para verificar que $H \leq G$, basta ver (1) que $h_{1}\ast h_{2} \in H  \quad \forall h_{1}, h_{2} \in H$, (2) que $1_{G} \in H$ y que (3), si $h \in H$ entonces $h^{-1} \in H$.

\end{namedstd}
\begin{namedres}{Teorema}{Caracterización de subgrupo}
Se $\emptyset \neq H \subseteq G$. Entonces $H \leq G$ si y solo si
\[
\forall x,y \in H  \quad (x \ast y^{-1} \in H).
\]
\end{namedres}
\begin{proof}
$(\implies)$ es trivial. Probemos $(\impliedby)$. Si $x \in H$, entonces $x \ast x ^{-1} = 1_{G}  \in H$. Por otro lado, como $1_{G} \in H$, $1_{G} \ast x^{-1} = x^{-1} \in H$. Finalmente, si $x,y \in H$, como $y^{-1} \in H$, entonces $x \ast (y ^{-1}) ^{-1} = x \ast y \in H$. Por tanto, $H$ es subgrupo.

\end{proof}
\begin{namedres}{Lema}{Intersección de subgrupos es subgrupo}
Sea $(G, \ast)$ y sea $\{ H_{i}:i \in I \}$ una familia de subgrupos de $G$. Entonces $H:=\bigcap_{i \in I} H_{i} \leq G_{i}$.

\end{namedres}
\begin{proof}
Note que $H\neq \emptyset$, pues $1_{G} \in H_{i}$ para todo $i$. Sea $x,y \in H$. Como $x,y \in H_{i}$ para todo $i \in I$. Entonces $x \ast y^{-1} \in H_{i}$ para todo $i$. Luego, $x \ast y ^{-1} \in H$.

\end{proof}
\begin{namedstd}{Definición}{Grupo generado por un subconjunto}
Sea $(G, \ast)$ un grupo y $S \subseteq G$. Defina
\[
\langle S \rangle := \bigcap \{ H:H\leq G  \quad \land  \quad S \subseteq H\}
\]
como el grupo generado por $S$.

\end{namedstd}
\begin{namedres}{Teorema}{Generado es el grupo más pequeño que contiene al conjunto}
$\langle S \rangle$ es el grupo más pequeño que contiene a $S$, i.e., si $K\leq G$ tal que $S \subseteq K$, entonces $\langle S \rangle \subseteq K$.
\end{namedres}
\begin{proof}
Si $K \leq G$ y $S \subseteq K$, entonces $K \in \{ H \leq G : S \subseteq H \}$, por lo que $\bigcap \{ H:H\leq G  \quad \land  \quad S \subseteq H\} = \langle S \rangle \subseteq K$.


\end{proof}
\begin{namedres}{Teorema}{Caracterización del grupo generado}
Si $(G, \ast)$ es un grupo y $S \neq \emptyset$, entonces
\[
\langle S \rangle = \{ x_{1}^{\alpha_{1}} \ast \dots \ast x_{n}^{\alpha_{n}}: n \in \mathbb{N}^{\ast}, x_{i} \in S, \alpha_{i} \in \{ -1,1\} \}.
\]
\end{namedres}
\begin{proof}
Sea $H = \{ x_{1}^{\alpha_{1}} \ast \dots \ast x_{n}^{\alpha_{n}}: n \in \mathbb{N}^{\ast}, x_{i} \in S, \alpha_{i} \in \{ -1,1\} \}$. Note que $H \neq \emptyset$. Sean $x,y \in H$. Entonces existen $k, \ell \in \mathbb{N}^{\ast}; \{ \alpha_{i} \}_{i=1}^{k}, \{ \beta_{i} \}_{i=1}^{\ell}$ con $\alpha_{i}, \beta_{i} \in \{ -1,1 \}$ para todo $i$, y $x_{1},\dots,x_{k},y_{1},\dots y_{\ell} \in S$ tales que
\[
x = x_{1}^{\alpha_{1}} \ast \dots \ast x_{k}^{\alpha_{k}}, y = y_{1}^{\beta_{1}}\dots y_{\ell}^{\beta_{\ell}}.
\]
Note que $x \ast y^{-1} = x_{1}^{\alpha_{1}} \ast \dots \ast x_{k}^{\alpha_{k}} \ast y_{1}^{-\beta_{1}}\dots y_{\ell}^{\beta_{-\ell}} \in H$. Así, $H \leq G$ y por tanto $\langle S \rangle \subseteq H$.
Por otro lado, como $\langle S \rangle$ es un grupo, $H \subseteq \langle S \rangle$.

\end{proof}
\begin{namedstd}{Definición}{Grupo generado}
SI $(G, \ast)$ es un grupo y $S \subseteq G$ tal que $\langle S \rangle = G$, decimos que $G$ es generado por $S$. Si $S$ es finito, decimos que $G$ es finitamente generado.

\end{namedstd}
\begin{namedstd}{Ejemplos}{de grupos generados}
\begin{enumerate}[label=(\arabic*)]
\item Sea $(G, \ast)$ un grupo. Entonces
\[
\begin{aligned}
\langle \emptyset \rangle &= \bigcap \{ H: H\leq  G \} \\
&= \{ 1_{G} \}.
\end{aligned}
\]
\item Si $(G, \ast)$ es un grupo y $S \leq G$, $\langle S \rangle = \langle S \rangle$.
\item Si $S = \{ i \} \subseteq (\mathbb{C}^{\ast}, \cdot)$, entonces $\langle S \rangle = \{ i,-1,1,-i \}$.
\item Si $S = \{ 2 \} \subseteq (\mathbb{Z},+)$, $\langle S \rangle = \{ 2k:k \in \mathbb{Z} \}$.
\item Si $S=\{ 12,-8 \} \subseteq (\mathbb{Z}, +)$, entonces $\langle S \rangle = \{ 12 a + (-8)b: a,b \in \mathbb{Z} \} = \{ 4k: k \in \mathbb{Z} \}$ .
\item En general, si $a \in (G, \ast)$,
\[
\begin{aligned}
\langle a \rangle &= \bigcap \{ H: H<a  \quad\land  \quad a \in H \}\\
&= \{ a^{n}: n \in \mathbb{Z} \}.
\end{aligned}
\]
\end{enumerate}

\end{namedstd}

\subsection{Grupos cíclicos}


\begin{namedstd}{Definición}{Orden de un grupo y de un elemento}
Sea $(G, \ast)$ un grupo. El orden de $G$ es la cantidad de elementos de $G$ y se denota por $\lvert  G \rvert$.
\begin{enumerate}[label=(\arabic*)]
\item Si existe $p$ primo tal que tal que $\lvert  G \rvert = p^{k}$ con $k \in \mathbb{N}$. decimos que $G$ es un $p$-grupo.
\item Si $G$ es un grupo y $g \in G$. el orden de $g$ ($\lvert  g \rvert$) es el menor $n \in \mathbb{N}$ tal que $g^{n} = 1_{G}$. Si tal $n$ no extiste, decimos que tal $\lvert g \rvert = \infty$.
\end{enumerate}

\end{namedstd}
\begin{namedstd}{Ejemplos}{Órdenes en $\Z_n$}
\begin{enumerate}[label=(\arabic*)]
\item $(\mathbb{Z}_{n}, +) = \{ [0],[1],\dots,[n-1] \}$, $\lvert \mathbb{Z}_{n} \rvert = n$.
\item Considere $(\mathbb{Z}_{4}, +) = \{ [0], [1], [2], [3] \}$, $\lvert \mathbb{Z}_{4} \rvert = 4$. Note que
\[
[2] \neq [0] \quad \land  \quad [2]+[2] = [0],
\]
\end{enumerate}
por lo que $\lvert [2] \rvert = 2$. De la misma manera, verificamos que $\lvert [1] \rvert = 4$ y $\langle [1] \rangle = \mathbb{Z}_{4}$.

\end{namedstd}
\begin{namedres}{Teorema}{Orden y divisibilidad}
Sea $G$ un grupo y $g \in G$, $n \in \mathbb{Z}$. Entonces
\[
g^{n} = 1_{G} \iff \lvert g \rvert \mid n.
\]
\end{namedres}
\begin{proof}
Hagamos el caso de $n \in \mathbb{N}$.
($\impliedby$): Si $\lvert g \rvert \mid n$, existe $k \in \mathbb{N}$ tal que $k \lvert g \rvert = n$. Entonces
\[
g^{n} = g^{k \lvert  g \rvert } = (g^{\lvert g \rvert })^{k} = (1_{G})^{k} = 1_{G}.
\]
$(\implies)$: Suponga que $g^{n} = 1_{G}$. Por minimalidad, $\lvert g \rvert \leq n$. Por algoritmo de la división, existen $q,r \in \mathbb{N}$ tales que $n = q \lvert g \rvert + r$ y $0 \leq r < \lvert g \rvert$. Entonces,
\[
1_{G} = g^{n} =g^{q \lvert g \rvert} \ast g^{r} = 1_{G} \ast g^{r} = g^{r},
\]
de donde concluimos que $r = 0$ por minimalidad de $\lvert g \rvert$, por tanto $\lvert g \rvert \mid n$.

Para el caso general $n \in \mathbb{Z}$, basta ver que $g^{n} = 1_{G} \iff g^{-n} = 1_{G}$ y además $\lvert  g \rvert \mid n \iff \lvert g \rvert \mid -n$.

\end{proof}
\begin{namedstd}{Definición}{Grupo cíclico}
Sea $G$ un grupo. $G$ es cíclico si existe $g \in G$ tal que $G = \langle g \rangle$.
\end{namedstd}
\begin{namedstd}{Ejemplo}{Grupos cíclicos}
\begin{enumerate}[label=(\arabic*)]
\item $(\mathbb{Z}, +) = \langle 1 \rangle = \langle -1 \rangle$. Tenemos que $\lvert \mathbb{Z} \rvert = \infty, \lvert 1 \rvert = \infty, \lvert 2 \rvert = \infty$.
\item Considere $\mathbb{Z}_{5}$ con la multiplicación de clases de equivalencia (congruencia módulo 5) definida anteriormente. Definimos
\[
U(5) = \{ [x] \in \mathbb{Z}_{5}: \exists [y] \in \mathbb{Z}_{5}([x] \cdot [y]=[1]) \} = \mathbb{Z}^{\ast}_{5}
\]
\end{enumerate}
¿$U(5)$ es un grupo cíclico? Sí, por ejemplo, $\lvert [2] \rvert = 4$. De la misma manera, su inverso $[3]$ también funciona para generar a $U(5)$. Pero $[4]$ no. Así, $\langle [2] \rangle, \langle [3] \rangle$.


\end{namedstd}
\begin{namedstd}{Definición}{Unidades de $\mathbb{Z}_{m}$}
Considere $\mathbb{Z}_{m}$. Defina
\[
U(m) = \{ [x] \in \mathbb{Z}_{m}: \exists [y] \in \mathbb{Z}_{m}([x] \cdot [y] = [1]) \}.
\]
Decimos que $[a]$ es una unidad de $\mathbb{Z}_{m}$ si $a \in U(m)$.

\end{namedstd}
\begin{namedstd}{Ejercicio}{Caracterización de unidades}
Muestre que $[a] \in U(m) \iff \MCD(a,m) = 1$.

\end{namedstd}
\begin{namedres}{Teorema}{Generador de $\Z_m$}
Sea $m \in \mathbb{N}^{\ast}$. Considere $(\mathbb{Z}_{m}, +)$. Sea $[a] \in \mathbb{Z}_{m}$. Sea $[a] \in \mathbb{Z}_{m}$. Entonces
\[
\langle [a] \rangle = \mathbb{Z}_{m} \iff \MCD(a,m) = 1.
\]
\end{namedres}
\begin{proof}
($\impliedby$): Suponga que $\MCD(a,m)=1$. Por el ejercicio anterior, existe $[a]^{-1} \in \mathbb{Z}_{m}$ tal que $[a] \cdot [a] ^{-1} = [1]$. Sea $[b] \in \mathbb{Z}_{m}$. Probaremos que $[b] \in \langle [a] \rangle$. Sea $k \in \mathbb{Z}$ tal que $[k]:=[b] \cdot [a^{-1}]$. Así.
\begin{align*}
[k]= [b] \cdot [a]^{-1} \implies [k] \cdot [a] = [b] \implies [a] + \dots + [a] = [b],
\end{align*}
de donde concluimos que $b \in \langle[a]\rangle$, y por tanto $\langle [a] \rangle = \mathbb{Z}_{m}$.
La otra dirección se deja como ejercicio.

\end{proof}
\begin{namedstd}{Nota}{Cardinalidad de grupos cíclicos}
Si $G$ es cíclico, entonces $G$ es numerable ($\lvert G \rvert \leq \lvert \mathbb{N} \rvert$). Por ser cíclico, existe $g \in G$ tal que $G = \langle g \rangle = \{ g^{n}:n \in \mathbb{Z} \}$. Considere la función $\Phi: \mathbb{Z} \to G$ tal que $n \mapsto g^{n}$. Como $g$ es cíclico, esta función es sobreyectiva. Así, $\lvert G \rvert \leq \lvert \mathbb{Z} \rvert$. Esto implica que cualquier grupo construido a partir de conjuntos no numerables, como $(\mathbb{R}, +)$, no puede ser cíclico.

\end{namedstd}
\begin{namedres}{Proposición}{Grupos cíclicos y orden}
Sea $G$ un grupo cíclico con $g \in G$ tal que $G = \langle g \rangle$. Entonces:
\begin{enumerate}[label=(\arabic*)]
\item Si $\lvert g \rvert = \infty$, entonces $g^{i}=g^{j} \iff i=j$.
\item Si $\lvert g \rvert=n \in \mathbb{N}$, entonces $\langle g \rangle = \{ e,g,\dots, g^{n-1}\}$. En este caso $g^{i} = g^{j} \iff n \mid i-j$. Esto implica que $\lvert G \rvert = \lvert g \rvert$.
\end{enumerate}
\end{namedres}
\begin{proof}
Para (1), suponga que $\lvert g \rvert=\infty$. Luego para todo $n \in \mathbb{N}$, $g^{n} \neq e$. Suponga que existen $i,j$ distintos (con $i<j$ sin pérdida de generalidad) tales que $g^{i} = g^{j}$. Como $j-i>0$, tenemos que $e = g^{i} \ast g^{-i} =g^{j-i}$, una contradicción pues $g$ tiene orden infinito.
Para (2), sabemos que $\langle g \rangle = \{ g^{n}: n \in \mathbb{Z} \}$, por lo que $\{ e,g,\dots, g^{n-1}\} \subseteq \langle g \rangle$. Ahora, sea $h \in \langle g \rangle$. Entonces, existe $k \in \mathbb{Z}$ tal que $h = g^{k}$. Por algoritmo de la división, existen $q,r \in \mathbb{Z}$ con $0\leq r\leq n-1$ tales que $k=nq+r$. Así, $h=g^{nq+r} = (g^{n})^{q} \ast g^{r}  = g^{r} \in \{ e,g,\dots, g^{n-1}\}$. Concluimos que $\langle  g \rangle = \{ e,g,\dots, g^{n-1}\}$.
Ahora, dados $i,j \in \mathbb{Z}$, suponga que $g^{i} = g^{j} \implies g^{i-j} = e$, luego $\lvert g \rvert \Big\lvert i-j$. Por último, basta mostrar que si $0 \leq i<j<n$ entonces $g^{i} \neq g^{j}$. Note que $0<j-i<n$, por lo que si $g^{i} = g^{j}$, entonces $g^{j-i} = e$, lo que contradice la minimalidad del orden $n$.

\end{proof}

\subsection{Subgrupos de Grupos Cíclicos}


\begin{namedres}{Lema}{Divisibilidad y subgrupo generado}
Sea $G$ un grupo cíclico y $g \in G$ tal que $G = \langle g \rangle$. Si $k \mid \ell$, entonces $\langle g^{\ell} \rangle \subseteq \langle g^{k} \rangle$.
\end{namedres}
\begin{proof}
Sea $x \in \langle g^{\ell} \rangle$. Existe $r \in \mathbb{Z}$ tal que $x=(g^{\ell})^{r}$. Además, como $k \mid \ell$, existe $t \in \mathbb{Z}$ tal que $tk = \ell$. Así, $x = (g^{tk})^{r} = (g^{k})^{tr} \in \langle g^{k} \rangle$.

\end{proof}
\begin{namedstd}{Ejercicio}{Orden de potencia}
Si $g \in G$ con $\lvert g \rvert = rs$, con $r,s \in \mathbb{N}$, entonces $\lvert g^{r} \rvert = s$.

\end{namedstd}
\begin{namedres}{Proposición}{Ordenes de subgrupos de grupos cíclicos}
Sea $G$ grupo y $g \in G$ tal que $\lvert g \rvert = n$. Dado $k \in \mathbb{N}$, entonces
\begin{enumerate}[label=(\arabic*)]
\item $\langle g^{k} \rangle = \langle g^{\mcd(n,k)} \rangle$
\item Si $c \mid n$, entonces $\lvert g^{c} \rvert  = \frac{n}{c}$.
\end{enumerate}
En particular, $\lvert g^{k} \rvert = \frac{n}{\mcd(n,k)}$.
\end{namedres}
\begin{proof}
Para (1), Sea $d=\mcd(n,k)$. Es claro que $d \mid k \implies \langle g^{k} \rangle \subseteq \langle g^{d} \rangle$. Por otro lado, por lema de Bezout, existen $r,t \in \mathbb{Z}$ tales que $d = rn + tk$. Luego, $g^{d} = g^{rn} \ast g^{tk} = (g^{k})^{t},$ por lo que $g^{d} \in \langle g^{k} \rangle$ y por tanto $\langle g^{d} \rangle \subseteq \langle g^{k} \rangle$. Concluimos que $\langle g^{d} \rangle = \langle g^{k} \rangle$.

Ahora, para (2), como $c \mid n$, existe $m \in \mathbb{Z}$ tal que $mc = n \implies m = \frac{n}{c} \in \mathbb{Z}$. Usando el ejercicio anterior, tenemos que $\lvert g \rvert = \frac{n}{c} \cdot c$ y por tanto $\lvert g^{c} \rvert = \frac{n}{c}$. Para el resultado particular, tome $c=d$ en (2) y use la propiedad (1).

\end{proof}
\begin{namedres}{Corolario}{orden de elementos en grupos cíclicos}
Sea $G$ cíclico tal que $\lvert G \rvert = n$ y $h \in G$. Entonces, $\lvert h \rvert \Big\lvert n$.
\end{namedres}
\begin{proof}
Sea $g \in G$ tal que $\langle g \rangle = G$. Como $h \in G$, existe $k \in \{ 0,\dots,n-1 \}$ tal que $h = g^{k}$. Así, $\lvert h \rvert = \lvert g^{k} \rvert = \frac{n}{\mcd(n,k)} \implies \lvert h \rvert \mcd(n,k) = n$, de donde se sigue la divisibilidad.

\end{proof}
\begin{namedres}{Teorema}{Igualdad de subgrupos}
Sea $G$ un grupo y $g \in G$ tal que $\lvert g \rvert = n$. Las siguientes son equivalentes:
\begin{enumerate}[label=(\arabic*)]
\item $\mcd(k,n) = \mcd(\ell, n)$
\item $\langle g^{k} \rangle = \langle g^{\ell} \rangle$.
\item $\lvert g^{k} \rvert = \lvert g^{\ell} \rvert$.
\end{enumerate}
\end{namedres}
\begin{proof}
$(2) \implies (3)$ es trivial. Para $(1) \implies (2)$, asuma que $\mcd(k,n) = \mcd(\ell, n)$. Note que
\[
\langle g^{k} \rangle = \langle g^{\mcd(n,k)} \rangle = \langle g^{\mcd(n,\ell)} \rangle = \langle g^{\ell} \rangle.
\]
Finalmente, para $(3) \implies (1)$, usando la proposición anterior, note que
\[
\lvert g^{k} \rvert = \lvert g^{\ell} \rvert \implies \frac{n}{\mcd(n,k)} = \frac{n}{\mcd(n, \ell)} \implies \mcd(n,k) = \mcd(n, \ell).
\]

\end{proof}
\begin{namedres}{Corolario}{Elemento generador de un grupo cíclico}
Sea $G = \langle g \rangle$ con $\lvert g \rvert=n$. Sea $k \in \mathbb{Z}$. Entonces $G = \langle g^{k} \rangle \iff \mcd(n,k) = 1$.
\end{namedres}
\begin{proof}
Por el teorema anterior $\langle g^{k} \rangle = \langle g^{1} \rangle \iff \mcd(n,k)=\mcd(n,1)  = 1$.

\end{proof}
\begin{namedres}{Teorema}{Subgrupo de grupo cíclico es cíclico}
Sea $G$ un grupo cíclico y $H \leq G$. Entonces $H$ es cíclico.
\end{namedres}
\begin{proof}
Sea $g \in G$ tal que $G = \langle g \rangle$. Si $H$ es el subgrupo trivial, claramente es cíclico. Suponga que $H$ es no trivial. Sea $e \neq h \in H$. Como $h \in H \subseteq G = \langle g \rangle$, existe $0 \neq k \in \mathbb{Z}$ tal que $h = g^{k}$. Defina $A = \{ n \in \mathbb{N}^{\ast} : g^{n} \in H\}$. Note que $A \neq \emptyset$, pues si $k \geq 0$, $k \in A$, y en caso contrario $-k \in A$ (por ser $H$ un grupo, $h^{-1} \in H$). Luego, aplicando el principio del buen orden, $A$ tiene un elemento minimal $\ell$.  Probaremos que $H = \langle g^{\ell} \rangle$. Por un lado, como $\ell \in A$, $g^{\ell}\in H$ y por tanto $\langle g^{\ell} \rangle \subseteq H$. Ahora, tome $t \in H$. Existe $i \in \mathbb{Z}$ tal que $t = g^{i}$. Suponiendo que $i \neq 0$, por algoritmo de la división, existen $q,r \in \mathbb{Z}$  con $0\leq r<\ell$ tales que $i = q\ell+r \implies r = i-q\ell$. Así, $g^{r} = g^{i} \ast g^{-q \ell} \in H$. Si $r\neq0$, tenemos que $r \in A$, lo que contradice la minimalidad de $\ell$. Concluya que $r=0$. Así, $t = g^{i} = (g^{\ell})^{q}$, de donde tenemos que $t \in \langle g^{\ell} \rangle$, Concluimos que $H = \langle g^{\ell} \rangle$.

\end{proof}
\begin{namedstd}{Nota}{Hallazgo del generador}
Por la prueba, sabemos como encontrar un generador de $H$ a partir de un generador de $G$,

\end{namedstd}
\begin{namedres}{Teorema}{Lagrange para grupos cíclicos}
Sea $G$ un grupo cíclico tal que $\lvert G \rvert = n$ y $H \leq G$. Entonces, $\lvert H \rvert \Big\lvert n$.
\end{namedres}
\begin{proof}
Sea $g \in G$ tal que $\langle g \rangle=G$ y $H\leq G$. Luego, $H$ es cíclico y existe $k \in \mathbb{N}$ tal que $H = \langle g^{k} \rangle$. Sea $d = \mcd(n,k)$. Entonces,
\[
\lvert H \rvert  = \lvert \langle g^{k} \rangle  \rvert = \lvert g^{k} \rvert = \frac{n}{d}  \implies d \lvert H \rvert = n,
\]
de donde se sigue la divisibilidad.

\end{proof}
\begin{namedres}{Teorema}{Unicidad de ordenes en subgrupos}
Sea $G$ un grupo cíclico con $\lvert G \rvert = n$ y sea $k \in \mathbb{N}$ tal que $k \mid n$. Entonces existe un único subgrupo $H$ de $G$ con $\lvert H \rvert = k$. Además, si $G = \langle g \rangle$ entonces $H = \langle g^{n/k} \rangle$
\end{namedres}
\begin{proof}
Sea $\ell = \frac{n}{k} \in \mathbb{N}$ (pues $k \mid n$). Sea $H = \langle g^{\ell} \rangle$. Entonces, $\lvert H \rvert = \lvert \langle g^{\ell} \rangle  \rvert = \frac{n}{\mcd(n, \ell)}$. Ahora, como $n = k\ell$ entonces $\mcd(n, \ell) = \mcd(k\ell, \ell) = \ell \mcd(k,1) = \ell$. Así, $\lvert H \rvert = \frac{n}{\ell} = k$. La unicidad queda como ejercicio.

\end{proof}

\subsection{Grupo de permutaciones}


Sea $n \in \mathbb{N}^{\ast}$. Defina
\[
S_{n} = \{ \sigma: \sigma \text{ es permutación de }\{ 1,\dots,n \} \}.
\]
Si $X$ es un conjunto cualquiera, $\sigma: X \to X$ biyectiva es una permutación.
$S_{X} = \{ \sigma : \quad \sigma: X \to X \text{ permutación} \}$. Es fácil ver que si $\lvert X \rvert = n$, $S_{X} \cong S_{n}$, pues
\[
\begin{aligned}
X &= \{ x_{1},\dots,x_{n} \}\\
&\cong \{ 1,\dots,n \} 
\end{aligned}.
\]

\begin{namedres}{Teorema}{Orden del grupo de permutaciones}
Sea $n \in \mathbb{N}^{\ast}$. Entonces $\lvert S_{n} \rvert = n!$.
\end{namedres}
\begin{proof}
Note que
\[
\begin{aligned}
\sigma(1) &\text{ tiene } n \text{ posibilidades}\\
\sigma(2) &\text{ tiene } n-1 \text{ posibilidades} \\
& \qquad \qquad \vdots \\
\sigma(n) &\text{ tiene } 1 \text{ posibilidad}.
\end{aligned}
\]

\end{proof}
\begin{namedstd}{Definición}{Mover y fijar un elemento}
Sea $\alpha \in S_{n}$ y sea $k \in \{ 1, \dots, n \}$. Decimos que:
\begin{enumerate}[label=(\arabic*)]
\item "$\alpha$ mueve $k$" si $\alpha(k) \neq k$,
\item "$\alpha$ fija $k$" si $\alpha(k) = k$.
\end{enumerate}

\end{namedstd}
\begin{namedstd}{Notación}{Composición de permutaciones}
Considere $\alpha:\{ 1,2,3\} \to \{ 1,2,3 \}$ y  $\beta:\{ 1,2,3\} \to \{ 1,2,3 \}$ tales que
\[
\alpha = \begin{pmatrix}
1 & 2 & 3 \\
2 & 1 & 3
\end{pmatrix}, \quad  
\beta = \begin{pmatrix}
1 & 2 & 3 \\
3 & 2 & 1
\end{pmatrix}.
\]
Tenemos que
\[
\alpha \circ \beta = \begin{pmatrix}
1 & 2 & 3 \\
2 & 1 & 3
\end{pmatrix} \circ \begin{pmatrix}
1 & 2 & 3 \\
3 & 2 & 1
\end{pmatrix} = \begin{pmatrix}
1 & 2 & 3 \\
3 & 1 & 2
\end{pmatrix},
\]

\[
\beta \circ \alpha = \begin{pmatrix}
1 & 2 & 3 \\
3 & 2 & 1
\end{pmatrix} \circ \begin{pmatrix}
1 & 2 & 3 \\
2 & 1 & 3
\end{pmatrix} = \begin{pmatrix}
1 & 2 & 3 \\
2 & 3 & 1
\end{pmatrix}.
\]
En general, dos permutaciones no tienen por qué conmutar.

\end{namedstd}
\begin{namedstd}{Definición}{r-ciclo}
Sea $n \in \mathbb{N}^\ast$ y $r \in \{ 1,\dots,n \}$. Decimos que $\alpha \in S_{n}$ es un $r$-ciclo. Si existen, $i_{1},\dots,i_{r} \in \{ 1,\dots,n \}$ distintos tales que
\[
\begin{aligned}
\alpha(i_{1}) &= i_{2}\\
\alpha(i_{2}) &= i_{3}\\
\vdots\\
\alpha(i_{r-1}) &= i_{r}\\
\alpha(i_{r}) &= i_{1}.
\end{aligned}
\]
Escribimos, $\begin{pmatrix}i_{1} & \cdots  & i_{r}\end{pmatrix}$. Note que si $\alpha = \begin{pmatrix} i_{1} & \cdots & i_{r}\end{pmatrix}$ y $j \not\in \{ i_{1},\dots,i_{r} \}$, entonces $\alpha(j) = j$.
\end{namedstd}
\begin{namedstd}{Ejemplo}{$r$-ciclo en $S_4$}
$\alpha=\begin{pmatrix}1 & 2 & 3 & 4 \\ 2 & 3 & 1 & 4\end{pmatrix} = \begin{pmatrix}1 & 2 & 3\end{pmatrix} \in S_{4}$.

\end{namedstd}
\begin{namedstd}{Definición}{Permutaciones disjuntas}
Sean $\alpha,\beta \in S_{n}$/ Decimos que $\alpha$ y $\beta$ son disjuntas si todo elemento que es movido por $\alpha$ es fijado por $\beta$ y todo elemento que mueve $\beta$ es fijado por $\alpha$. Note que esta definición admite que haya elementos fijados por ambas,

\end{namedstd}
\begin{namedstd}{Definición}{Transposición}
Decimos que un $2$-ciclo es una transposición.
\end{namedstd}
\begin{namedstd}{Nota}{Notación cíclica}

\[
\begin{pmatrix}
1 & 3 & 5
\end{pmatrix} = \begin{pmatrix}
5 & 1 & 3
\end{pmatrix} = \begin{pmatrix}
3 & 5 & 1
\end{pmatrix}
\]
son el mismo ciclo.

\end{namedstd}
\begin{namedstd}{Ejercicio}{$r$-ciclo y orden}
Muestre que $\alpha \in S_{n}$ es un $r$-ciclo si y solo si $\lvert a \rvert = r$ (i.e., $\alpha^{r} = \mathrm{id}_{n}$).

\end{namedstd}
\begin{namedstd}{Nota}{Inverso de un ciclo}
Si $\alpha = \begin{pmatrix}1 & 3 & 5\end{pmatrix} \in S_{5}$, $\alpha ^{-1} = \begin{pmatrix}5 & 3 & 1\end{pmatrix}$.

\end{namedstd}
\begin{namedstd}{Ejemplo}{Producto de transposiciones disjuntas}
Considere $\alpha \in S_{4}$ con
\[
\alpha = \begin{pmatrix}
1 & 2 & 3 & 4 \\
2 & 1 & 4 & 3
\end{pmatrix} = \underbrace{ \begin{pmatrix}
1 & 2
\end{pmatrix} }_{ \beta_{1} } \underbrace{ \begin{pmatrix}
3 & 4
\end{pmatrix} }_{ \beta_{2} },
\]
donde $\beta_{1}, \beta_{2} \in S_{1}$ son $2$-ciclos disjuntos.

\end{namedstd}
\begin{namedres}{Teorema}{Conmutatividad de permutaciones disjuntas}
Sean $\alpha, \beta \in S_{n}$ disjuntas. Entonces $\alpha \circ \beta = \beta \circ \alpha$.
\end{namedres}
\begin{proof}
Sea $i \in \{ 1,\dots, n \}$. Tenemos tres casos posibles.
\textbf{Caso 1:} $\alpha$ mueve $i$. Sea $j = \alpha(i) \neq i$. Entonces, $\beta(i) = i$. Note además que $\alpha(j) \neq j$, pues de lo contrario, $\alpha(i) = j = \alpha(j)$ implica por inyectividad de $\alpha$ que $i = j$, una contradicción. Por tanto, $\beta(j) = j$. Así, tenemos que
\[
\begin{aligned}
(\alpha \circ \beta)(i) &= \alpha(\beta(i)) = \alpha(i) = j \\ 
(\beta \circ \alpha)(i) &= \beta(\alpha(i)) = \beta(j) = j.
\end{aligned}
\]
\textbf{Caso 2:} $\beta$ mueve $i$. Este caso es análogo al anterior.
\textbf{Caso 3:} $\alpha, \beta$ fijan a $i$: Note que $\alpha(i) = \beta(i) = i$. Luego
\[
\begin{aligned}
(\alpha \circ \beta)(i) &= \alpha(\beta(i)) = \alpha(i) = i \\ 
(\beta \circ \alpha)(i) &= \beta(\alpha(i)) = \beta(j) = i.
\end{aligned}
\]
Concluya que permutaciones disjuntas conmutan.

\end{proof}
\begin{namedres}{Teorema}{Orbita finita}
Sea $\alpha \in S_{n}$, $i \in \{ 1,\dots,n \}$ tal que $\alpha(i) \neq i$. Defina la sucesión $\{ i_{k} \}_{k \in \mathbb{N}^{\ast}}$ de manera que $i_{1} = i$, $i_{2} = \alpha(i_{1}) = \alpha(i), \dots,i_{k+1} = \alpha (i_{k})$ para todo $k \in \mathbb{N}^{\ast}$. Entonces:
\begin{enumerate}[label=(\arabic*)]
\item existe $k \in \{ 1,\dots,n \}$ tal que $i_{k+1} \in \{ i_{1},\dots, i_{k} \}$;
\item si $r = \min \{ k \in \mathbb{N}^{\ast} : i_{k+1} \in \{ i_{1},\dots,i_{k} \} \} \implies i_{r+1} = i_{1}.$
\end{enumerate}
\end{namedres}
\begin{proof}
Para (1), note que para todo $k \in \mathbb{N}^{\ast}$, $\{ i_{1},\dots,i_{k} \} \subseteq \{ 1,\dots,n \}$ y $\{ 1,\dots,n \}$ es finito, entonces el resultado se sigue directamente.
Para (2), defina $r$ como en el enunciado. Tome $i_{j}:= i_{r+1} \in \{ i_{1},\dots,i_{r} \}$. Suponga por contradicción que $j>1 \implies j-1>0$. Note que $\alpha^{j-1}(i) = i_{j} = i_{r+1} = \alpha^{r}(i)$. Aplicando $\alpha ^{-1}$ $j-1$ veces, tenemos que
\[
i_{1} = i = \alpha^{r+1-j}(i),
\]
pero $r+1-j < r$, por lo que esto contradice la minimalidad de $r$. Por tanto $j = 1$.

\end{proof}
\begin{namedres}{Lema}{Igualdad de ciclos}
Sean $\alpha, \beta \in S_{n}$ y sea $i \in \{ 1,\dots,n \}$. Si $\alpha, \beta$ mueven $i$ y para todo $k \in \mathbb{N}$, $\alpha^{k}(i) = \beta^{k}(i)$, entonces $\alpha = \beta$.

\end{namedres}
\begin{proof}
ejercicio.
\end{proof}
\begin{namedres}{Teorema}{Factorización en ciclos disjuntos}
Toda permutación en $S_{n}$ es producto de ciclos disjuntos. Además, si agregamos los $1$-ciclos la descomposición es única módulo el orden.

\end{namedres}
\begin{proof}
Sea $\alpha \in S_{n}$. Proceda por inducción fuerte sobre el número $k$ de elementos que mueve $\alpha$. Si $k=0$, entonces $\alpha = \mathrm{id}_{n}$. Entonces, $\alpha = (1)(2) \cdots (n)$.
Considere ahora el caso de $k>0$. Suponga que el teorema es cierto para $\{0,1,\dots,k-1\}$. Sea $i \in \{ 1,\dots,n \}$ tal que $\alpha(i) \neq i$. Defina $i_{k} \in \{ 1,\dots,n \}$ tal que $i_{1} = i$ y $i_{m+1} = \alpha(i_{m}) = \alpha^{m}(i)$ para $m \in \mathbb{N}^{\ast}$. Sea $r = \min \{ \ell \in \mathbb{N}: i_{\ell+1} \in \{ i_{1},\dots,i_{\ell} \} \} \leq n$. Por teorema anterior, $i_{r+1}=i_{1}$. Sea $\delta = \begin{pmatrix}i_{1} & \cdots  & i_{r}\end{pmatrix} \in S_{n}$ un $r$=ciclo.
Si $r = n$ , $\delta$ es un $n$-ciclo y $\delta = \alpha$. Si $r<n$, defina $A := \{ 1,\dots,n \} \setminus \{ i_{1},\dots,i_{r} \}$. Note que $\alpha(A) = A$, pues todos estos puntos son fijados por $\alpha$. Además, $\delta$ fija todos los elementos de $A$ por construcción. Sea $\beta \in S_{n}$ tal que para todo $x \in A$, $\beta(x) = \alpha(x)$ y para todo $x \in \{ i_{1},\dots,i_{r} \}$, $\beta(x) = x$. Note que $\beta$ y $\delta$ son disjuntas. Además, $\alpha = \beta \circ \delta$. El número de elementos que $\beta$ mueve es menor que $k$. Luego, por hipótesis inductiva aplicada a $\beta$, $\beta$ es producto de ciclos disjuntos, y por tanto $\alpha = \beta \circ \delta$ también lo es.

\end{proof}
\begin{namedres}{Proposición}{Descomposición en transposiciones}
Sea $n\geq 2$. Todo $\alpha \in S_{n}$ es producto de transposiciones.
\end{namedres}
\begin{proof}
Basta mostrarlo si $\alpha$ es un ciclo. Sea $\alpha = \begin{pmatrix}a_{1} & \cdots & a_{n}\end{pmatrix}$ un ciclo. Proceda por inducción fuerte sobre $r$. Para $r=1$, note que $\alpha = (a)$. Luego, $\alpha$ es la identidad de $S_n$. Sea $b \in \{ 1,\dots,n \}$ tal que $b\neq a$  (sabemos que existe porque $n>2$). Luego, $\alpha = \begin{pmatrix} a & b\end{pmatrix} \begin{pmatrix}a & b\end{pmatrix} = \mathrm{id}_{S_{n}}$. Suponga como hipótesis inductiva que el teorema es cierto para cada $k$-ciclo con $k \leq r-1$.
Probaremos que $\alpha = \begin{pmatrix}a_{1} & a_{3} & \cdots  & a_{r-1}\end{pmatrix} \begin{pmatrix}a_{1} & a_{2}\end{pmatrix}$. Sea $\beta = \begin{pmatrix}a_{1} & a_{3} & \cdots  & a_{r}\end{pmatrix} \begin{pmatrix}a_{1} & a_{2}\end{pmatrix}$. Note que
\[
\begin{aligned}
\beta(a_{r}) = a_{1} &= \alpha(a_{r}) \\
\beta(a_{1}) = a_{2} &=\alpha(a_{1})\\
\beta(a_{2}) = a_{3} &= \alpha(a_{2})\\
\beta(a_{k}) = a_{k+1} &= \alpha(a_{k}), \qquad 3\leq k\leq r-1.\\b
\beta(b) = b &= \alpha(b), \qquad b \not\in \{ a_{1},\dots,a_{k} \}.
\end{aligned}
\]
Luego $\alpha = \beta$. Aplicamos H.I. en $\begin{pmatrix}a_{1} & a_{3} & \cdots a_{r}\end{pmatrix}$.

\end{proof}
\begin{namedstd}{Nota}{Forma estándar}
$\alpha = \begin{pmatrix} x_{1} & \cdots & x_{m}\end{pmatrix} = \begin{pmatrix}x_{1} & x_{m}\end{pmatrix} \begin{pmatrix}x_{1} & x_{m-1}\end{pmatrix} \cdots \begin{pmatrix}x_{1} & x_{2}\end{pmatrix}$.

\end{namedstd}
\begin{namedstd}{Ejemplo}{Factorización en transposiciones}
Considere
\[
\begin{aligned}
\alpha &=\begin{pmatrix}
1 & 2 & 3 & 4 & 5 & 6 & 7 & 8 & 9 \\ 
6 & 2 & 7 & 9 & 8 & 3 & 1 & 5 & 4 \\
\end{pmatrix} \\
&= \begin{pmatrix}
1 & 6 & 3 & 7
\end{pmatrix} \begin{pmatrix}
2
\end{pmatrix} \begin{pmatrix}
4 & 9
\end{pmatrix} \begin{pmatrix}
8 & 5
\end{pmatrix} \qquad \text{factorización en ciclos} \\
&= \begin{pmatrix}
1 & 3 & 7
\end{pmatrix} \begin{pmatrix}
1 & 6
\end{pmatrix} \begin{pmatrix}
1 & 2
\end{pmatrix} \begin{pmatrix}
1 & 2
\end{pmatrix} \begin{pmatrix}
4 & 9
\end{pmatrix} \begin{pmatrix}
5 & 8
\end{pmatrix} \\
&=\begin{pmatrix}
1 & 7
\end{pmatrix} \begin{pmatrix}
1 & 3
\end{pmatrix} \begin{pmatrix}
1 & 6
\end{pmatrix} \begin{pmatrix}
1 & 2
\end{pmatrix} \begin{pmatrix}
1 & 2
\end{pmatrix} \begin{pmatrix}
4 & 9
\end{pmatrix} \begin{pmatrix}
5 & 8
\end{pmatrix} \qquad \text{descomposición en transposiciones}.
\end{aligned}
\]


\end{namedstd}
\begin{namedstd}{Definición}{Signo}
Sea $n \in \mathbb{N}$, $\alpha \in S_{n}$. Si $\alpha = \beta_{1} \cdots \beta_{k}$ es una factorización completa en ciclos disjuntos. Defina $\sgn(\sigma) = (-1)^{n-k}$.

\end{namedstd}
\begin{namedres}{Teorema}{Signo y transposiciones}
Sea $\alpha \in S_{n}$ y $\tau$ una transposición. Entonces, $\sgn(\tau\alpha) = -\sgn(\alpha)$.
\end{namedres}
\begin{proof}
Ejercicio

\end{proof}
\begin{namedres}{Teorema}{Signo separa producto}
Sea $\alpha, \beta \in S_{n}$. Entonces $\sgn(\alpha\beta) = \sgn(\alpha) \sgn(\beta)$.
\end{namedres}
\begin{proof}
Ejercicio, inducción sobre el número de transposiciones.

\end{proof}
\begin{namedstd}{Definición}{Paridad de la permutación}
Decimos que $\alpha \in S_{n}$ es par si $\alpha$ se puede descomponer como el producto de un número par de transposiciones. En caso contrario $\alpha$ es impar.

\end{namedstd}
\begin{namedres}{Teorema}{Paridad y signo}
Sea $\alpha \in S_{n}$. $\alpha$ es impar si y solo si $\sgn( \alpha ) = - 1$.

\end{namedres}
\begin{namedstd}{Definición}{Grupo alternante}
$A_{n}:= \{ \sigma \in S_{n}: \sigma \text{ es par} \}$ es un grupo. Además. $\lvert A_{n} \rvert = \frac{n!}{2}$.

\end{namedstd}

\subsection{Centro y centralizador}


\begin{namedstd}{Definición}{centro de $G$}
Sea $(G, \ast)$ un grupo. Defina
\[
Z(G) = \{ g \in G: \forall h \in G (gh = hg)\} 
\]
\end{namedstd}
\begin{namedstd}{Nota}{Propiedades del centro}
\begin{enumerate}[label=(\arabic*)]
\item $e \in Z(G)$,
\item $G$ es abeliano si y solo si  $Z(G) = G$.
\item Si $g \in Z(G)$, entonces para todo $h \in G$, $h = g \ast h \ast g^{-1}$.
\end{enumerate}


\end{namedstd}
\begin{namedres}{Teorema}{Centro es subgrupo}
Si $G$ es un grupo, entonces $Z(G) \leq G$ y $Z(G)$ es conmutativo.

\end{namedres}
\begin{namedstd}{Definición}{Centralizador}
Sea $G$ un grupo y $a \in G$. Defina el centralizador como
\[
C(a):= \{ g \in G: a \ast g = g \ast a \}
\]

\end{namedstd}
\begin{namedstd}{Ejercicio}{Centralizador es subgrupo}
$C(a) \leq G$.

Se puede probar que $Z(G) = \bigcap_{a \in G} C(a)$.
\end{namedstd}

% ======================================================================
\section{Homomorfismos de Grupos y Subgrupos Normales}
% ======================================================================

\subsection{Definiciones elementales}

\begin{namedstd}{Definición}{Homomorfismo de grupos}
Sean $(G,\ast_{G}), (H,\ast_{H})$ grupos. Un homomorfismo de grupos entre $G$ y $H$ es una función $f:G\to H$ tal que
\[
\forall a,b, \in G  \quad \Big( f(a \ast_{G} b) = f(a) \ast_{h} f(b)\Big).
\]

\end{namedstd}
\begin{namedstd}{Ejemplos}{Homomorfismos básicos}
\begin{enumerate}[label=(\arabic*)]
\item Sea $G$ un grupo y $f:G\to G$ tal que $f(x) = 1_{G}$ para todo $x \in G$. Note que
\[
f(x \ast y) = 1_{G} = 1_{G} \ast 1_{G} = f(x) \ast f(g)
\]
\item Considere $f:(\mathbb{R},+) \to (\mathbb{R}^{\ast}, \cdot)$ tal que $f(x) = e^{x}$. Note que
\[
f(x+y) = e^{x+y} = e^{x}
 \cdot e^{y} = f(x) \cdot f(y).
\]
\end{enumerate}

\end{namedstd}
\begin{namedres}{Proposición}{Propiedades de los homomorfismos}
Sean $(G,\ast_{G}), (H,\ast_{H})$ grupos y $f:G\to H$ un homomorfismo. Entonces:
\begin{enumerate}[label=(\arabic*)]
\item $f(1_{G}) = 1_{H}$,
\item $\forall g \in G  \quad \Big( f(g^{-1}) = \big(f(g)\big)^{-1}\Big)$,
\item $\forall n \in \mathbb{Z} \quad \Big(f(g^{n}) = \big(f(g)\big)^{n}\Big)$
\item Si $g \in G$ y $\lvert g \rvert = n$, entonces $\lvert f(g) \rvert \Big\lvert n$.
\end{enumerate}
\end{namedres}
\begin{proof}
Para (1), note que
\[
f(1_{G}) = f(1_{G} \ast_{G} 1_{G}) = f(1_{G}) \ast_{H} f(1_{G}),
\]
de donde concluimos que $f(1_{G}) = 1_{H}$.
Para (2), note que
\[
g \ast_{G} g^{-1} = 1_{G} \implies f(g) \ast_{H} f(g^{-1}) = 1_{H} \implies f(g^{-1}) = \big(f(g)\big)^{-1}.
\]
Para (3), en el caso $n\geq 0$, procedemos por inducción. El caso base $n=0$ se probó en (1). Suponga como hipótesis inductiva que $f(g^{n})=\big(f(g)\big)^{n}$ para algún $n \in \mathbb{N}$. Note que
\[
\begin{aligned}
f(g^{n+1}) = f(g^{n} \ast_{G} g) &=f(g^{n}) \ast_{H} f(g)\\
&=\big(f(g)\big)^{n} \ast_{H} f(g) \\
&= \big(f(g)\big)^{n+1}.
\end{aligned}
\]
Ahora, para $n<0$, note que $-n > 0$. Así, $g^{n} = (g^{-n})^{-1}$. Así,
\[
f(g^{n}) = f\big((g^{-n})^{-1}\big) = \big(f(g^{-n})\big)^{-1} = \Big(\big(f(g)\big)^{-n}\Big)^{-1} = \big(f(g)\big)^{n}.
\]
Finalmente, para (4), como $\lvert g \rvert = n$,
\[
g^{n} = 1_{G} \implies f(g^{n}) = \big(f(g)\big)^{n} = 1_{H},
\]
de donde deducimos que $\lvert f(g) \rvert \Big\lvert n$.

\end{proof}
\begin{namedstd}{Ejemplo}{Proyección y homomorfismos a $\Z$}
\begin{enumerate}[label=(\arabic*)]
\item Considere $f:(\mathbb{Z}, +)\to(\mathbb{Z}_{n}, +)$ tal que $x \mapsto [x]$. Note que
\[
f(x+y) = [x+y] = [x] + [y] = f(x) + f(y),
\]
\end{enumerate}
de donde concluimos que $f$ es un homomorfismo.
\begin{enumerate}[label=(\arabic*)]
\item ¿Existe un homomorfismo no trivial $f:\mathbb{Z}_{n} \to \mathbb{Z}$? Si existiera, note que $\lvert [1] \rvert = n$ y entonces $\lvert f[1] \rvert \Big\lvert n$, una contradicción, pues el orden de los elementos en $\mathbb{Z}$ es infinito (salvo el cero).
\end{enumerate}

\end{namedstd}
\begin{namedres}{Proposición}{Composición de homomorfismos}
Sean $G,H,K$ grupos y $f:G\to H$, $g:H\to K$ homomorfismos. Entonces $g \circ f:G \to K$ es un homomorfismo.
\end{namedres}
\begin{proof}
Ejercicio

\end{proof}
\begin{namedstd}{Definición}{Kernel e Imagen de un homomorfismo}
Sean $G,H$ grupos y $f:G\to H$ homomorfismo. Defina
\[
\begin{aligned}
\Ker(f) &= \{ x \in G: f(x) = 1_{H} \} \subseteq G \\
\Imm(f) &= \{ f(x):x \in G \} \subseteq H.
\end{aligned}
\]

\end{namedstd}
\begin{namedres}{Teorema}{Kernel e Imagen son subgrupos}
Sean $G,H$ grupos y $f:G\to H$ homomorfismo. Entonces,
\begin{enumerate}[label=(\arabic*)]
\item $\Ker(f) \leq G$
\item $\Imm(f) \leq H$
\end{enumerate}

\end{namedres}
\begin{proof}
Para (1), note que $1_{G} \in \Ker(f)$, por lo que $\Ker(f) \neq \emptyset$. Sean $x,y \in \Ker(f)$.
\[
f(x \ast_{G} y^{-1}) = f(x) \ast_{H} f(y^{-1}) =1_{H} \implies x \ast_{G} y \in \Ker(f).
\]
Para (2), note que $f(1_{G}) = 1_{H} \in \Imm(f)$, por lo que $\Imm(f) \neq \emptyset$. Sean $a,b \in \Imm(f)$. Luego, existen $x,y \in G$ tales que $f(x) = a$ y $f(y) = b$. Por propiedades de homomorfismos, $b^{-1} = f(y^{-1})$. Así,
\[
f(x \ast_{G} y^{-1}) = f(x) \ast_{H} \big(f(y)\big)^{-1} = a \ast_{H} b^{-1} \in \Imm(f).
\]

\end{proof}
\begin{namedstd}{Ejemplo}{Proyección $\Z \to \Z_n$}
Sea $\pi:\mathbb{Z} \to \mathbb{Z}_{n}$ tal que $x \mapsto [x]$. Tenemos que
\[
\begin{aligned}
\Imm(\pi) &= \{ \pi(x):x \in \mathbb{Z} \} \\
&=\{ [x]: x \in \mathbb{Z} \} \\
&=\mathbb{Z}_{n},\\
\Ker(\pi) &=\{ x:\pi(x) = [0] \} \\
&=\{ x: [x] = 0 \} \\
&=n \cdot \mathbb{Z}.
\end{aligned}
\]

\end{namedstd}
\begin{namedstd}{Definición}{Monomorfismos, epimorfismos e isomorfismos}
Sea $f:G\to H$ un homomorfismo de grupos:
\begin{enumerate}[label=(\arabic*)]
\item si $f$ es inyectivo, decimos que $f$ es un monomorfismo,
\item si $f$ es sobreyectivo, decimos que $f$ es un epimorfismo,
\item si $f$ es biyectivo, decimos que $f$ es un isomorfismo.
\end{enumerate}

\end{namedstd}
\begin{namedres}{Teorema}{Existencia del isomorfismo inverso}
Si $f:G\to H$ es un isomorfismo, entonces $f^{-1}:H\to G$ (como función) es un isomorfismo.
\end{namedres}
\begin{proof}
Ejercicio.

\end{proof}
\begin{namedres}{Teorema}{Monomorfismos y kernel}
Sea $f:G\to H$ homomorfismo. Entonces,
\[
f \text{ es monomorfismo} \iff \Ker(f) = \{ 1_{G} \}
\]
\end{namedres}
\begin{proof}
ejercicio. Notar que $f(x)=f(y) \implies f(x \ast y^{-1}) = 1_{H}$.


\end{proof}
\begin{namedres}{Teorema}{Imagen inversa y homomorfismos}
Sea $f:G\to H$ un homomorfismo y sea $H'\leq H$. Entonces $f^{-1}(H') \leq G$
\end{namedres}
\begin{proof}
Note que $1_{G} \in f^{-1}(H')$ pues $f(1_{G}) = 1_{H} \in H \implies f^{-1}(H') \neq \emptyset$. Sean $x,y \in f^{-1}(H')$. Hay que mostrar que $x \ast_{G} y^{-1} \in f^{-1}(H')$. Como $f(x) \in H', f(y) \in H'$, entonces $\big(f(y)\big)^{-1} = f(y^{-1}) \in H'$. Entonces $f(x \ast_{G} y^{-1}) = f(x)\ast_{H} f(y^{-1}) \in H'$.

\end{proof}
\begin{namedstd}{Definición}{Automorfismo y grupo de automorfismos}
Sea $G$ un grupo. Un automorfismo de $G$ es un isomorfismo de $G$ en $G$. Defina $\mathrm{Aut} \hspace{2pt}(G) = \{ f:G\to G \text{ isomorfismo}\}$ como el conjuto de automorfismos de $G$.
Note que $\mathrm{Aut} \hspace{2pt}(G) \subseteq S_{G} = \{ \phi: G\to G \text{ biyectiva}\}$.

\end{namedstd}
\begin{namedres}{Teorema}{Subgrupo de automorfismos}
$\mathrm{Aut} \hspace{2pt}(G) \leq S_{g}$.
\end{namedres}
\begin{proof}
Ejercicio

\end{proof}
\begin{namedstd}{Definición}{Grupos isomorfos}
Decimos que dos grupos $G$ y $H$ son isomorfos si existe un isomorfismo $f:G\to H$. Escribimos $G \cong H$.

\end{namedstd}
\begin{namedstd}{Ejemplo}{Isomorfismo $\Z_4 \cong G$}

Considere $G = (\{ 1,i,-1,-i \}, \ast)$.
\begin{center}\fbox{\itshape Figura del cuaderno (20260407\_144838.jpg) omitida}\end{center}
\begin{center}\fbox{\itshape Figura del cuaderno (20260407\_144841.jpg) omitida}\end{center}

$Z_{4} \cong G$ bajo el isomorfismo
\[
\begin{aligned}
0  &\mapsto 1 \\
1 &\mapsto i \\
2 &\mapsto -1 \\
3 &\mapsto -i
\end{aligned}
\]

\end{namedstd}
\begin{namedstd}{Nota}{Ser isomorfo es relación de equivalencia}
Note que la relación $G \cong H$ es reflexiva, transitiva y simétrica (es fácil probarlo). Luego es una relación de equivalencia.

\end{namedstd}
\begin{namedres}{Teorema}{Isomorfismos de grupos cíclicos}
Si $G = \langle g \rangle$, entonces:
\begin{enumerate}[label=(\arabic*)]
\item Si $\lvert g \rvert = \infty \implies G \cong \mathbb{Z}$.
\item Si $\lvert g \rvert = n \implies G \cong \mathbb{Z}_{n}$.
\end{enumerate}

\end{namedres}

\subsection{Cadenas y sucesiones exactas}


\begin{namedstd}{Definición}{Complejo de cadenas}
Un complejo de cadenas es un conjunto $\mathcal{S} = (G_{i}, \phi_{i})_{i \in \mathbb{Z}}$ en donde:
\begin{enumerate}[label=(\arabic*)]
\item Para todo $i \in \mathbb{Z}$, $G_{i}$ es un grupo,
\item Para todo $i \in \mathbb{Z}$, $\phi_{i}: G_{i} \to G_{i-1}$ es un homomorfismo de grupos.
\item $\phi_i \circ \phi_{i+1}: G_{i+1} \to G_{i-1}$ es tal que $(\phi_{i} \circ \phi_{i+1})(x) = 1_{G_{i-1}}$ para todo $x \in G_{i+1}$.
\end{enumerate}

\[
\cdots \to G_{3} \overset{\phi_{3}}{\to} G_{2} \overset{\phi_{2}}{\to} G_{1} \overset{\phi_{1}}{\to} G_{0} \overset{\phi_{-1}}{\to} G_{-1} \overset{\phi_{-2}}{\to} \cdots
\]
Note que $(\phi_{i} \circ \phi_{i+1})(x) = \phi_{i} \big(\phi_{i+1}(x)\big)= 1_{G_{i-1}}$, luego $\phi_{i+1}(x) \in \Ker(\phi_{i})$ para todo $i \in \mathbb{N}$. Sea $Z_{n}(\mathcal{S}) = \Ker(\phi_{n}) \subseteq G_{n}$. $B_{n}(\mathcal{S})=\Imm(\phi_{n+1}) \subseteq G_{n}$.  Note que
\[
B_{n}(\mathcal{S}) \leq  Z_{n}(\mathcal{S}) \leq  G_{n}
\]

\end{namedstd}
\begin{namedstd}{Definición}{Sucesión exacta}
Tome  $\mathcal{S} = (G_{i}, \phi_{i})_{i \in \mathbb{Z}}$ un complejo de cadenas. Si tenemos que para todo $i$ tenemos que $\Imm(\phi_{i+1}) = \Ker(\phi_{i})$. Decimos que $\mathcal{S}$ es una sucesión exacta.

\end{namedstd}
\begin{namedstd}{Definición}{Sucesión corta}
Una sucesión es corta si es finita.

\end{namedstd}
\begin{namedstd}{Ejemplo}{Sucesión exacta corta}
Considere
\[
0 \rightarrow G_{2} \overset{ f }{\rightarrow} G_{1} \overset{ g }{\rightarrow} G_{0} \rightarrow 0,
\]
donde $0$ denota el grupo trivial. Los homomorfismos sin nombre son los triviales.

\end{namedstd}
\begin{namedres}{Teorema}{Exactitud de la sucesión e inyectividad/sobreyectividad}
Sea $\phi_{0}:G_{0}\to G_{1}$ homomorfismo:
\begin{enumerate}[label=(\arabic*)]
\item Si $0 \overset{ f }{\rightarrow} G_{0} \overset{ \phi_{0} }{\rightarrow} G_{1}$ es exacta, entonces $\phi_{0}$ es inyectiva.
\item Si $G_{0} \overset{ \phi_{0} }{\rightarrow} G_{1} \overset{ f }{\rightarrow} 0$ es exacta, entonces $\phi_{0}$ es sobreyectiva.
\item Si $0 \overset{  f }{\rightarrow} G_{0} \overset{ \phi_{0} }{\rightarrow} G_{1} \overset{ g }{\rightarrow} 0$, entonces $\phi_{0}$ es un isomorfismo.
\end{enumerate}
\end{namedres}
\begin{proof}
Para (1), note que $\Ker(\phi_{0}) = \Imm(f) = \{ 1_{G_{0}} \}$, por lo que $\phi_{0}$ es inyectiva. Para (2), note que $\Imm(\phi_{0}) = \Ker f = G_{1}$. Combinamos ambos resultados para (3).

\end{proof}
\begin{namedstd}{Ejemplo}{Sucesión exacta $\Z \to \Z_2$}
Considere
\[
0 \rightarrow 2\mathbb{Z} \overset{ i }{\rightarrow}  \mathbb{Z} \overset{ \pi }{\rightarrow} \mathbb{Z}_{2} \rightarrow 0,
\]
donde $i(x) = x$ y $\pi(x) = [x]$ para todo $x \in \mathbb{Z}$. Basta mostrar que $\Imm(i) = \Ker(\pi)$.

\end{namedstd}

\subsection{Traslación y conjugación}


\begin{namedstd}{Definición}{Traslación}
Sea $(G, \ast)$ un grupo y sea $a \in G$. La traslación izquierda (dercha) de $a$ es $T_{a}:G\to G$ tal que $x \mapsto a \ast x$ ($x \mapsto x \ast a$).

\end{namedstd}
\begin{namedstd}{Nota}{Traslación es biyectiva}
\begin{enumerate}[label=(\arabic*)]
\item $T_{a}$ es inyectiva pues $T_{a}(x)=T_{a}(y) \iff a \ast x = a \ast y \iff x = y$.
\item $T_{a}$ es sobreyectiva, pues si $y \in G$, $T_{a}(a^{-1} \ast y) = y$.
\end{enumerate}
Luego, $T_{a}$ es biyectiva. Concluya que $T_{a} \in S_{G} = \{ \sigma:G\to G \text{ biyectiva}\}$.

\end{namedstd}
\begin{namedres}{Teorema}{Traslaciones de $G$ son un monomorfismo}
Sea $f:G\to S_{G}$ tal que $f(a) = T_{a}$. Entonces $f$ es un monomorfismo.
\end{namedres}
\begin{proof}
Probaremos primero que $f$ es un homomorfismo, i.e., que $f(a \ast b) = T_{a \ast b} = T_{a} \ast T_{b}$. Sea $x \in G$. Tenemos que
\[
\begin{aligned}
(T_{a} \circ T_{b})(x) = T_{a}(T_{b}(x)) &= T_{a} (b \ast x)\\
&= a \ast (b \ast x) \\
&=(a \ast b) \ast x \\
&= T_{a \ast b} (x)
\end{aligned},
\]
de donde tenemos que $f$ es un homomorfismo. Probaremos que $f$ es inyectiva. Sea $a \in \Ker(f)$. Sea $a \in  G$ tal que $f(a) = T_{a} = \mathrm{id}_{G}$. Luego,
\[
\forall x \in G  \quad(T_{a}(x)= x) \implies \forall x \in G  \quad(a \ast x = x) \implies a = 1_{G}.
\]
Así, $\Ker(f) = \{ \mathrm{id_{G}} \}$. Concluimos que $f$ es un monomorfismo.

\end{proof}
\begin{namedstd}{Definición}{Conjugación}
Sea $G$ un grupo y $a \in G$. Defina la conjugación por $a$ $C_{a}: G\to G$ tal que $x \mapsto a \ast x \ast a^{-1}$.

\end{namedstd}
\begin{namedres}{Teorema}{Conjugación en el automorfismo}
Sea $G$ un grupo y $a \in G$. Entonces $C_{a} \in \mathrm{Aut}(G)$.
\end{namedres}
\begin{proof}
Note que $C_{a}(x \ast y) = C_{a}(x) \ast C_{a}(y)$ (ejercicio), por lo que $C_{a}$ es homomorfismo. Para probar que es inyectiva, tome $x \in \Ker(C_{a})$. Luego
\[
C_{a}(x) = 1_{G} \implies a \ast x \ast a ^{-1} = 1_{G} \implies x = 1_{G}
\]
Para probar que es sobreyectiva, sea $x \in G$. Note que $C_{a}(a^{-1} \ast x \ast a) = x$, de donde concluimos sobreyectividad.

\end{proof}
\begin{namedres}{Teorema}{Conjugaciones de $G$ son un homomorfismo}
Sea $f:G\to \mathrm{Aut}(G) \leq S_{G}$ tal que $f(a) = C_{a}$. Entonces $f$ es un homomorfismo.
\end{namedres}
\begin{proof}
Sea $x \in G$. Entonces
\[
\begin{aligned}
(C_{a} \circ C_{b})(x) &= C_{a} (C_{b}(x))\\
&=C_{a} ( b \ast x \ast b^{-1})\\
&=a \ast ( b \ast x \ast b^{-1}) \ast a ^{-1} \\
&=(a \ast b) \ast x \ast (a \ast b) ^{-1} \\
&=C_{a \ast b}(x).
\end{aligned}
\]

\end{proof}

\subsection{Coclases}


\begin{namedstd}{Definición}{Coclases}
Sea $G$ un grupo y $H \leq G$:
\begin{enumerate}[label=(\arabic*)]
\item Una coclase izquierda (derecha) de $H$ es $aH = \{ a \ast h:h \in H \}$ ($Ha = \{ h \ast a:h \in H \}$), con $a \in G$.
\end{enumerate}
Se suele escribir "$a+H$" o "$H+a$" en grupos aditivos.

\end{namedstd}
\begin{namedstd}{Ejemplos}{Coclases en $\Z$ y $S_3$}
\begin{enumerate}[label=(\arabic*)]
\item Considere $G = (\mathbb{Z}, +)$ y sea $n \in \mathbb{Z}$. Tome $H = \langle n \rangle = n \mathbb{Z}$. Sea $a \in \mathbb{Z}$. Note que
\[
\begin{aligned} 
a + H &= \{ a+h:h \in H \} \\
&=\{ a + nk: k \in \mathbb{Z} \} \\
&=\{ t: t \equiv a  \quad \mathrm{mod } \hspace{2mm}n \}.
\end{aligned}
\]
\item Considere $G = S_{3}$ y $\tau = \begin{pmatrix}1 & 2\end{pmatrix}$. Tome $H = \langle \tau \rangle = \{ \mathrm{id}, \tau \}$. Tenemos que
\[
\begin{aligned}
H \cdot \begin{pmatrix}
1 & 2 & 3
\end{pmatrix} &= \{ \begin{pmatrix}
1 & 2 & 3
\end{pmatrix}, \begin{pmatrix}
1 & 2
\end{pmatrix} \begin{pmatrix}
1 & 2 & 3
\end{pmatrix} \} \\
&=\{ \begin{pmatrix}
1 & 2 & 3
\end{pmatrix}, \begin{pmatrix}
2 & 3 
\end{pmatrix} \},\\
\begin{pmatrix}
1 & 2 & 3
\end{pmatrix} \cdot H &= \{ \begin{pmatrix}
1 & 2 & 3
\end{pmatrix}, \begin{pmatrix}
3 & 1
\end{pmatrix} \},\\
H \cdot \begin{pmatrix}
3 & 2 & 1
\end{pmatrix} &= \{ \begin{pmatrix}
3 & 2 & 1
\end{pmatrix}, \begin{pmatrix}
1 & 3
\end{pmatrix} \}
\end{aligned}
\]
\end{enumerate}
Note que $H \cdot \begin{pmatrix}3 & 2 & 1\end{pmatrix} \quad \cap \quad H \ast \begin{pmatrix}1 & 2 & 3\end{pmatrix} = \emptyset$.

\end{namedstd}
\begin{namedres}{Lema}{Criterio de igualdad de coclases}
Sea $H \leq G$ y $a,b \in G$. Entonces
\[
aH = bH \iff (b^{-1} \ast a) \in H.
\]
\end{namedres}
\begin{proof}
($\implies$): Tenemos que $a \in bH$. Luego, existe $h \in H$ tal que $a = b \ast h \implies b^{-1} \ast a = h \in H$.
($\impliedby$): Suponga que $b ^{-1} \ast a \in H$. Luego, existe $h_{0} \in H$ tal que $b ^{-1} \ast a = h_{0}$. Entonces $a = b \ast h_{0}$, entonces $a \in bH$. Luego, $aH \subseteq bH$, pues si $x \in aH$, tenemos que
\[
x = a \ast h = b \ast\underbrace{ h_{0} \ast h }_{ \in H }\in bH
\]
La otra inclusión se prueba de manera análoga.

\end{proof}
\begin{namedres}{Teorema}{Dos cualesquiera coclases son iguales o disjuntas}
Sea $H\leq G$ y sean $a,b \in G$. Entonces $aH = bH$ o $aH \cap bH = \emptyset$.
\end{namedres}
\begin{proof}
Suponga que $aH \cap bH \neq \emptyset$. Sea $x \in aH \cap bH$. Luego, existen $h_{1}, h_{2} \in H$ tal que $x= a \ast h_{1} = b \ast h_{2}$. Entonces, $b^{-1} \ast x = b^{-1} \ast a \ast h_{1} \in H$ , de modo que $b ^{-1} \ast a \in H$. Por el lema anterior $aH = bH$., de donde concluimos el resultado.

\end{proof}
\begin{namedstd}{Nota}{Coclases como partición}
Esto nos muestra que las coclases forman una partición del grupo. Más aún, si definimos la relación $\sim$ en $G$ como $a \sim b \iff aH = bH$, esta es una relación de equivalencia.

\end{namedstd}
\begin{namedstd}{Definición}{Conjunto cociente}
Definimos el conjunto cociente de $G$ sobre $H$  como $G/H := \{ aH:a \in G \}$, esto es, el conjunto de representantes de $G$ cocientado con la clase de equivalencia de coclases de $H$.

\end{namedstd}
\begin{namedres}{Teorema}{Igualdad del número de coclases izquierdas/derechas}
Sea $H\leq G$. Entonces el número de coclases derechas es igual al número de coclases izquierdas.
\end{namedres}
\begin{proof}
Sean $A = \{ aH: a \in G \}$ y $B = \{ Hb:b \in G \}$. Sea $f:A\to B$ tal que $f(aH) = Ha^{-1}$. $f$ está bien definida: Si $aH = a'H$, entonces $H a^{-1} = H(a')^{-1}$. Terminar la prueba (ejercicio).

\end{proof}
\begin{namedres}{Proposición}{Cardinalidad de subgrupo y coclases}
Si $H\leq G$ y $a \in G$, entonces $\lvert a H \rvert = \lvert H \rvert$ (como cardinalidad pues $aH$ no es necesariamente un grupo)
\end{namedres}
\begin{proof}
Defina $f:H\to aH$ tal que $h \mapsto a \ast h$. Note que
\[
\begin{aligned}
f(h_{1}) = f(h_{2}) \implies a \ast h_{1} = a \ast h_{2} \implies a^{-1} \ast a \ast h_{1} = a^{-1} \ast a \ast h_{2} \implies h_{1}=h_{2}.
\end{aligned}
\]

\end{proof}
\begin{namedstd}{Definición}{Índices}
Sea $H \leq G$. Definimos el índice de $H$ en $G$ como $[G:H] = \lvert G/H \rvert$ es el número de coclases diferentes de $H$ en $G$.

\end{namedstd}
\begin{namedres}{Teorema}{Lagrange para grupos finitos}
Sea $G$ un grupo finito y $H\leq G$. Entonces, $\lvert G \rvert = [G:H] \cdot \lvert H \rvert$. En particular, $\lvert H \rvert \Big\lvert \lvert G \rvert$.
\end{namedres}
\begin{proof}
Sea $k = [G :H]$. Sean $a_{1}H, a_{2}H, \dots, a_{k}H$ las coclases diferentes de $H$ en $G$. Entonces $G = \bigcup_{i=1}^{k} a_{i} H$. Como $a_{1}H, a_{2}H, \dots, a_{k}H$ son disjuntos, entonces $\lvert G \rvert = \sum_{i=1}^{k} \lvert a_{i} H \rvert = k \lvert H \rvert$.

\end{proof}
\begin{namedres}{Corolario}{Orden de un elemento divide al orden del grupo}
Sea $a \in G$. Entonces $\lvert a \rvert \Big\lvert \lvert G \rvert$.
\end{namedres}
\begin{proof}
Tome $H = \langle a \rangle$. Entonces $\lvert H \rvert = \lvert a \rvert$ y el resultado se sigue del teorema de Lagrange.

\end{proof}
\begin{namedres}{Corolario}{Multiplicatividad del índice}
Sean $K,H,G$ grupos tales que $K \leq H\leq G$. Entonces, $[G:K] = [G:H] [H:K]$.
\end{namedres}
\begin{proof}
Aplicando el teorema de Lagrange sobre los tres índices:
\[
[G:H] \cdot \lvert H \rvert = \lvert G \rvert, \quad [H:K] \cdot \lvert K \rvert = \lvert H \rvert, \quad [G:H] \cdot [H : K] = \lvert G \rvert. 
\]
Así,
\[
[G:H] \cdot [H:K] = \frac{\lvert G \rvert}{\lvert H \rvert } \cdot \frac{\lvert H \rvert}{\lvert K \rvert }  = \frac{\lvert G \rvert}{\lvert K \rvert } = [G:K].
\]

\end{proof}
\begin{namedres}{Proposición}{Coclases sobre el kernel de un homomorfismo}
Sean $G, G'$ y $f:G \to G'$ un homomorfismo. Sea $H = \Ker(f) \leq G$. Sea $a \in G$ y $a' = f(a)$. Entonces $aH = \{ x \in G:f(x) = a' \}$.
\end{namedres}
\begin{proof}
Sea $x \in aH$, entonces existe $h \in H$ tal que $x = a \ast h$. Luego, $f(x) = f(a) \ast_{G'} f(h) = f(a) = a'$. Luego, $aH \subseteq \{ x \in G:f(x) = a' \}$. Por otro lado, si $x \in G$ tal que $f(x) = a'$, tenemos que
\[
f(x) = f(a) \iff \big(f(a)\big)^{-1} \ast_{G'} f(x) = 1_{G'} \iff f(a^{-1} \ast x) = 1_{G'},
\]
de donde tenemos que $a^{-1} \ast x \in H$ y, por el criterio de igualdad de coclases, $xH = aH$, de donde $x \in aH$. Así, $\{ x \in G: f(x) = a' \} \subseteq aH$, y por tanto $aH \subseteq \{ x \in G:f(x) = a' \}$.

\end{proof}

\subsection{Subgrupos normales}


\begin{namedstd}{Definición}{Producto de conjuntos}
Sea $G$ un grupo y $S,T \subseteq G$. Defina $ST = \{ s \ast t: s \in S, t \in T \}$.
Note que si $S = \{ s \}$ y $T\leq G$, entonces $sT =ST$, una coclase de $G$. Note que si $S,T,U \subseteq G$, . Si $\mathcal{F} = \{ A:A\subseteq G \}$, entonces $\cdot$ está definida en $\mathcal{F}$ y es asociativa.

\end{namedstd}
\begin{namedstd}{Definición}{Subgrupo normal}
Sea $H \leq G$. Decimos que $H$ es un subgrupo normal de $G$ y escribimos $H \triangleleft G$ si $Hx = xH$ para todo $x \in G$.

\end{namedstd}
\begin{namedstd}{Nota}{Propiedades de subgrupos normales}
\begin{enumerate}[label=(\arabic*)]
\item Si $G$ es abeliano, todo subgrupo es normal.
\item $xH = Hx \iff xH x^{-1} = \{ x \ast h \ast x ^{-1}: h \in H \} = H$.
\item Si $H \triangleleft G$, esto no implica que para todos $x \in G$ y $h \in H$ $x \ast h = h \ast x$. Lo que sí es cierto es que si $h \in H$, existe $h' \in H$ tal que $x \ast h = h' \ast x$.
\item Si para todo $x \in  G$, $x H x ^{-1} \subseteq H$, esto necesariamente implica la igualdad entre ambos conjuntos, y por tanto, que $H \triangleleft G$. Suponga que $xHx ^{-1} \subseteq H$ para todo $x \in G$. Note que, para $x \in G$.
\[
H = x ^{-1}(\underbrace{  x H x ^{-1} }_{ \in H }) x \subseteq x ^{-1} H x \subseteq H \implies x ^{-1} H x = H.
\]
\item Para todo grupo $G$, $\{ 1_{G} \} \triangleleft G$, $G \triangleleft G$.
\end{enumerate}

\end{namedstd}
\begin{namedres}{Proposición}{Subgrupos encadenados}
Sean $K,H,G$ grupos con $K\leq H\leq G$. Si $K \triangleleft G$, entonces $K \triangleleft H$.
\end{namedres}
\begin{proof}
El resultado es trivial, pues $G \subseteq H$.


\end{proof}
\begin{namedres}{Teorema}{Kernel es subgrupo normal}
Sea $f:G\to G'$ un homomorfismo de grupos. Entonces, $\Ker(f) \triangleleft G$.
\end{namedres}
\begin{proof}
Sea $H:= \Ker(f)$. Sea $x \in G$. Hay que mostrar que $xHx ^{-1} \subseteq H$. Sea $h \in H$. Entonces, $f(x \ast h \ast x ^{-1}) = f(x) \ast f(h) \ast f(x ^{-1}) = 1_{G'}$.

\end{proof}
\begin{namedstd}{Ejemplo}{Subgrupo normal en $S_3$}
Sea $G = S_{3} = \{ (1), \begin{pmatrix}1 & 2\end{pmatrix}, \begin{pmatrix}1 & 3\end{pmatrix}, \begin{pmatrix}2 & 3\end{pmatrix} \begin{pmatrix}1 & 2 & 3\end{pmatrix} \begin{pmatrix}3 & 2 & 1\end{pmatrix} \}$. Considere, $H = \langle \begin{pmatrix}1 & 2 & 3\end{pmatrix} \rangle$.  Tenemos que
\begin{center}\fbox{\itshape Figura del cuaderno (20260417\_150411.jpg) omitida}\end{center}
y así por el resto de ciclos.
\end{namedstd}

\subsection{Grupo cociente}


\begin{namedres}{Proposición}{Grupos normales y operaciones}
Si $H \triangleleft G$ y $a,b \in G$, entonces $aH \cdot bH = (a \ast b) H$
\end{namedres}
\begin{proof}
Note que
\[
\begin{aligned}
aH \cdot bH
\end{aligned} = aH \cdot Hb = a Hb = a b H = (a \ast b) H.
\]

\end{proof}
\begin{namedres}{Teorema}{Existencia del grupo cociente}
Sea $H \triangleleft G$. Entonces $G / H$ es un grupo bajo la operación $\cdot$
\[
a H \cdot bH = (a \ast b) H.
\]
\end{namedres}
\begin{proof}
La operación está bien definida y es asociativa (véase proposición anterior).
\begin{enumerate}[label=(\arabic*)]
\item Pertenencia de la identidad: $1_{G} \cdot H = H$
\item Existencia de inversos: $aH \cdot a^{-1}H = (a \ast a^{-1}) H = H = id_{G / H }$.
\end{enumerate}

\end{proof}
\begin{namedstd}{Nota}{Independencia del representante}
¿Por qué no depende del representante? Suponga que $aH = \hat{a}H$ y $bH = \hat{b}H$. Entonces
\[
aH \cdot bH = \hat{a} H \cdot \hat{b}H \iff (a \ast b) H = (\hat{a} \ast \hat{b}) H \iff (a \ast b) \ast (\hat{ a} \ast \hat{b})^{-1} \in H,
\]
que es claramente cierto.

\end{namedstd}
\begin{namedres}{Teorema}{Todo subgrupo normal es kernel de un homomorfismo}
Sea $H \triangleleft G$ y considere $f:G\to G / H$ tal que $f(a) = a H$. Entonces $f$ es un epimorfismo y $\Ker(f) = H$.
\end{namedres}
\begin{proof}
Probaremos primero que $f$ es un homomorfismo. Note que $f(a \ast b) = (a \ast b) H = aH \cdot bH = f(a) \cdot f(b)$. La sobreyectividad se sigue por la definición. Ahora, si $x \in H$, entonces $f(x) = xH = H =id_{G / H}$, luego, $x \in \Ker f$. Si $x \in \Ker f$, entonces $f(x) = xH = H$, luego $x \in H$.

\end{proof}
\begin{namedstd}{Ejemplo}{Cocientes triviales}
\begin{enumerate}[label=(\arabic*)]
\item Note que $\{ 1_{G} \} \triangleleft G \implies G / \{ 1_{G} \} = \{ a \cdot \{ 1_{G}: a \in G \} \} = \{ \{ a \}: a \in G \} \cong G$.
\item $G \triangleleft G \implies G / G = \{ a \cdot G: a \in G \} = \{ G \} = \{ 1_{G / G} \}$.
\end{enumerate}

\end{namedstd}
\begin{namedstd}{Ejemplo}{Cocientes en $\Z$ y en grupo cíclico}
\begin{enumerate}[label=(\arabic*)]
\item Sea $G = (\mathbb{Z}, + )$, $H = 4\mathbb{Z}$. Note que $H \triangleleft G$ pues $G$ es abeliano. Note que $G / H \cong \mathbb{Z}_{4}$, pues $aH = bH \iff a + (-b) \in 4\mathbb{Z} \iff 4 \mid a-b \iff a \cong b  \quad \text{mod }4$.
\item Sea $G$ un grupo y sea $a \in G$ tal que $\lvert a \rvert = 8$. y $G=\langle a \rangle$. Sea $H = \langle a^{4}. \rangle$ Como $G$ es conmutativo. $H \triangleleft G$. Tenemos que $H = \{ 1_{G, }a^{4} \}$ y $G = \{ 1_{G}, a, a^{2},\dots,a^{7} \}$. ¿Quién es $G/H$? Encontremos las coclases:
\begin{enumerate}[label=(\alph*)]
\item $H = a^{4}H$
\item $aH = \{ a, a^{5} \} = a^{5} H$
\item $a^{2}H = \{ a^{2}, a^{6} \} = a^{6}H$
\item $a^{3}H = \{ a^{3},a^{7} \}= a^{7}H$
Así, $G/H = \{ H, aH, a^{2}H, a^{3}H \}$ y $G = H  \hspace{2mm}  \dot{\cup} \hspace{2mm} aH \hspace{2mm} \dot{\cup} \hspace{2mm} a^{2}H \hspace{2mm}\dot{\cup} \hspace{2mm} a^{3}H$.
\end{enumerate}
\end{enumerate}


\end{namedstd}
\begin{namedstd}{Definición}{Conmutador y subgrupo conmutador}
Sea $(G,\ast)$ un grupo. Sean $a,b \in G$. Defina el conmutador de $a,b$ como
\[
[a,b] = a \ast b \ast a^{-1} \ast b ^{-1} \in G
\]
Defina el subgrupo conmutador de $G$ como $G':=  \langle [a,b]:a,b \in G \rangle$.

\end{namedstd}
\begin{namedstd}{Ejercicio}{Caracterización via conjugación}
Sea $S < G$. Entonces $S \triangleleft G$ si y solo si  $C_{a}(S) \leq S$ para todo $a \in G$.
\end{namedstd}
\begin{namedres}{Teorema}{Subgrupo conmutador y cocientes abelianos}
Sea $G$ un grupo y $G'$ su subgrupo conmutador. Entonces $G' \triangleleft G$. Además, si $H \triangleleft G$, entonces $G/H$ es abeliano si y solo si $G' \triangleleft H$.
\end{namedres}
\begin{proof}
Hay que probar que, para todo $g \in G$, $gG'g^{-1} \subseteq G'$. Basta mostrar que si $a, b \in G$, $g \in G$ , entonces $g[a,b]g^{-1} \in G'$, ¿por qué? Si $z = [a_{1} , b_{1}] [a_{2}, b_{2}]$, entonces
\[
gzg^{-1} = g [a_{1},b_{1}] [a_{2}, b_{2}] g^{-1} = (g[a_{1},b_{1}]g^{-1})(g [a_{2},b_{2}] g^{-1}),
\]
luego, basta probar el resultado para cualquier conmutador y se sigue inmediatamente para todo $z \in G'$, Sea $[a,b] \in G'$. Entonces,
\[
\begin{aligned}
g \ast a \ast b \ast a^{-1} \ast b^{-1} \ast g^{-1}
\end{aligned}
\]
\[
= (g \ast a \ast g^{-1}) \ast (g \ast b \ast g^{-1}) \ast(g \ast a^{-1} \ast g ^{-1}) \ast (g \ast b ^{-1} \ast g ^{-1}) = [g \ast a \ast g ^{-1},g \ast b \ast g ^{-1} ]\in G'.
\]
Suponga ahora que $G/H$ es abeliano. Entonces para todos $a,b \in G$, $aH \cdot b H = b H \cdot a H$, por lo que $(a \ast b) H = (b \ast a)H$. Luego , $(a \ast b) \ast (b \ast a) ^{-1} = a \ast b \ast a^{-1} \ast b ^{-1} = [a,b] \in H$. Así, $\{ [a,b]: a,b \in G \} \subseteq H \implies G' \subseteq H \implies G' \leq H$. (Escribir la otra dirección).


\end{proof}
\begin{namedstd}{Ejemplo}{$A_3 \triangleleft S_3$}
Considere $S_{3} = \{ (1),\begin{pmatrix}1 & 2\end{pmatrix}, \begin{pmatrix}2 & 3\end{pmatrix}, \begin{pmatrix}1 & 3\end{pmatrix}, \begin{pmatrix}1 & 2 & 3\end{pmatrix} \begin{pmatrix}3 & 2 & 1\end{pmatrix} \}$ y $A_{3} = \{ \sigma \in S_{3}: \sigma \text{ es par} \} = \{ \begin{pmatrix}1\end{pmatrix}, \begin{pmatrix}1 & 2 & 3\end{pmatrix}, \begin{pmatrix}3 & 2 & 1\end{pmatrix}\}$. Probaremos que $A_{3} \triangleleft S_{3}$. Considere $S = \{ -1,1 \}$. Note que $(S, \cdot) = (\mathbb{Z}_{2}, +)$. Sea $f:S_{3} \to S$ tal que $\sigma \mapsto \sgn( \sigma )$. Note que $f$ es un homomorfismo, pues
\[
\begin{aligned}
f(\sigma_{1} \circ \sigma_{2}) &= \sgn(\sigma_{1} \circ \sigma_{2}  )\\
&= \sgn( \sigma_{1} ) \cdot \sgn( \sigma_{2} )\\
&= f(\sigma_{1}) \cdot f(\sigma_{2}).
\end{aligned}
\]
Note que $\Ker f = A_{3}$, de donde concluimos que $A_{3} \triangleleft S_{3}$. ¿Quién es $S_{3}/A_{3}$? Construyamos las coclases
\begin{enumerate}[label=(\arabic*)]
\item $A_{3}$
\item $\begin{pmatrix}1 & 2\end{pmatrix} A_{3} = \{ \begin{pmatrix}1 & 2\end{pmatrix}, \begin{pmatrix}3 & 2\end{pmatrix}, \begin{pmatrix}1 & 3\end{pmatrix} \}$.
\end{enumerate}
Así, $S_{3}/A_{3} = \{ A_{3}, \begin{pmatrix}1&2\end{pmatrix}A_{3} \} \cong \mathbb{Z}_{2}$.

\end{namedstd}
\begin{namedstd}{Ejemplo}{$\langle r_1\rangle \triangleleft D_4$}
Considere $D_{4} = \{ r_{0},r_{1},r_{2},r_{3},s_{1},s_{2},s_{3},s_{4} \}$ el grupo de simetrías del cuadrado y $H = \langle r_{1} \rangle = \{ r_{0},r_{1},r_{2},r_{3} \}$. Debemos probar que para todo $i \in \{ 1,\dots,4 \}$, $s_{i} H s_{i} ^{-1} \subseteq H$. Note en primer lugar que para todo $i$, $s_{i}^{2} = r_{0}$, luego, $s_{i} = s_{i}^{-1}$.  Luego, hay que probar que para todos $j \in \{ 0,1,2,3 \}, i \in\{ 1,2,3,4 \}$, $s_{i} r_{j} s_{i} ^{-1} \in H$


\end{namedstd}

\subsection{Teoremas del Isomorfismo}


\begin{namedres}{Teorema}{Primer teorema del isomorfismo}
Sean $G, G'$ grupos y sea $f:G \to G'$ un homomorfismo y sea $H = \Ker(f)$. Entonces, $G / H \cong \Imm(f)$.
\end{namedres}
\begin{proof}
Como $H \triangleleft G$, $G / H$ es un grupo. Sea $f':G / H \to G'$ tal que $aH \mapsto f(a)$. Probaremos que $f'$ está bien definida, i.e., que envía miembros de la misma clase de equivalencia a las mismas imágenes. Dados $a,b \in G$, suponga que $aH = bH$, i.e., que $b^{-1}a \in H$. Así,
\[
f(b^{-1}a)=1_{G'} \implies (f(b))^{-1} f(a) = 1_{G'}\implies f(b) = f(a) \implies f'(aH) = f'(bH).
\]
Probaremos ahora que $f'$ es un homomorfismo. Dados, $aH, bH \in G / H$, note que
\[
f'(aH \cdot bH) = f'((a \ast b)H)= f(a \ast b) =f(a) \ast_{G'}f(b) = f'(aH)\ast_{G'} f'(bH).
\]
Así, $f$ es homomorfismo. Para la inyectividad, sea $aH \in \Ker(f')$. Luego,
\[
f'(aH) = 1_{G'} \implies f(a) = 1_{G'} \implies a \in \Ker(f) = H \implies aH=H.
\]
Finalmente, probaremos que $f'$ es sobreyectiva, i.e, que $\Imm(f') = \Imm(f)$. Note que
\[
b \in \Imm(f') \iff \exists a \in G (f'(aH)=b) \iff \exists a \in G (f(a)=b) \iff b \in \Imm(f).
\]
Así, $f'$ es un isomorfismo, de donde concluimos el resultado.

\end{proof}
\begin{namedres}{Teorema}{Tercer teorema del isomorfismo}
Sean $K \leq H \leq G$ grupos tales que $K \triangleleft G$ y $H \triangleleft G$. Entonces $H / K \triangleleft G / K$ y $G / H \cong (G/K) / (H / K)$.
\end{namedres}
\begin{proof}
Probaremos primero que $H / K \triangleleft G / K$. Sea $aK \in G / K$. Hay que mostrar que $(aK) \cdot H /K \cdot (aK)^{-1} \subseteq H / K$. Sea $bK \in H / K$, luego
\[
(aK) \cdot (bK) \cdot (aK)^{-1} = (aba^{-1})K \in H / K.
\]
Concluimos que $H / K \triangleleft G / K$. Ahora, defina $f:G/K \to G / H$ tal que $aK \mapsto aH$. Probaremos que $f$ está bien definida. Tome $a,b \in G$ tales $aK = bK$. Entonces,
\[
b^{-1}a \in K \subseteq H \iff b^{-1} a \in H \iff aH = bH \implies f(aK) = f(bK).
\]
Probemos que $f$ es un homomorfismo:
\[
f(aK \cdot bK) = f(a \cdot b \cdot K) \quad= ab H = aH \cdot bH =f(aK) \cdot f(bK),
\]
de donde concluimos que $f$ es un homomorfismo. Finalmente, note que
\[
aK \in\Ker(f) \iff f(aK) =H \iff aH =H \iff a \in H \iff aK \in H / K,
\]
de modo que $\Ker(f)  = H / K$. Finalmente, aplicando el primer teorema del isomorfismo, concluimos el resultado.

\end{proof}
\begin{namedres}{Teorema}{conmutatividad de subgrupos normales}
Sea $N,H \leq G$ con $N \triangleleft G$. Entonces $NH=HN$.
\end{namedres}
\begin{proof}
Ejercicio

\end{proof}
\begin{namedres}{Teorema}{Segundo teorema del isomorfismo}
Sea $(G, \cdot)$ un grupo, $H \leq G$, $N \triangleleft G$. Entonces $N \cap H \triangleleft H$, $NH \leq G$ y $H / (N \cap H) \cong NH / N$.
\end{namedres}
\begin{proof}
Sabemos que $N \cap H \subseteq H$, $N \cap H \subseteq N$. Sea $x \in H$. Entonces $x \in G$. Hay que mostrar que $x(N \cap H)x ^{-1} \subseteq N \cap H$. Evidentemente, $x(N \cap H)x ^{-1} \subseteq H$. Como $N \triangleleft G$, entonces $x(N \cap H)x ^{-1} \subseteq x N x ^{-1} \subseteq N$. Así, $x(N \cap H)x ^{-1} \subseteq N \cap H$ y probamos normalidad de $N \cap H$.
La prueba de $NH\leq G$ se deja como ejercicio.
Finalmente, defina $f:H \to G / N$ tal que $h \mapsto h N$. Note que $f$ es un homomorfismo, pues $f(g \cdot h) = g h N = (gN)(hN) = f(g) \cdot f(h)$. Además,
\[
h \in \Ker(f) \iff f(h) = N \iff hN =N \iff h \in N,
\]
por lo que $\Ker(f) = N \cap H$. Finalmente, probamos que $\Imm(f) = NH /N$. Sea $gN \in \Imm(f)$. Luego, existe $h \in H$ tal que $hN = gN$. Luego, $g \in HN = NH$, entonces $gN \in NH / N$. Ahora, sea $g \in HN / N$. Entonces, existe $x \in HN$ tal que $g = xN \in HN / N$. Luego, $x = h_{1} n_{1}$, con $h_{1} \in H$, $n_{1} \in N$. Así,
\[
g = xN = (h_{1} \cdot n_{1})N = h_{1} N,
\]
entonces $g = f(h_{1}) \in \Imm(f)$. El resultado se sigue al aplicar el primer teorema del isomorfismo.

\end{proof}
\begin{namedstd}{Ejercicio}{Imagen inversa de un subgrupo}
Si $f:G \to G^{\ast}$ es un homomorfismo y $S^{\ast}\leq G^{\ast}$. Entonces $f^{-1}(S^{\ast}) = \{ x \in G  : f(x)\in S^{\ast} \}$

\end{namedstd}
\begin{namedres}{Teorema}{Correspondencia}
Sea $(G, \cdot)$ un grupo, $H \triangleleft G$ y $p:G \to G / H$ tal que $p(g)=gH$. Entonces $K \mapsto p(K) = K / H$ es una biyección entre los subgrupos de $G$ que contienen a $H$ y los subgrupos de $G/H$. Además, si $K$ es tal que $H\leq K\leq G$ y $K^{\ast}:=K / H$, tenemos que
\begin{enumerate}[label=(\arabic*)]
\item $L\leq K \iff L^{\ast} \leq K^{\ast}$ y en este caso $[K:L] =[K^{\ast}:L^{\ast}]$.
\item $L \triangleleft K \iff L^{\ast} \triangleleft K^{\ast}$, y en este caso, $K / L \cong K^{\ast} / L^{\ast}$.
\end{enumerate}

\end{namedres}
\begin{proof}
Sean $A = \{ L: H\leq L \leq G \}$ y $B = \{ \tilde{G}: \tilde{G} \leq G / H \}$. Defina $\Phi:A \to B$ tal que $\Phi(L) = L / H$. Suponga que $L,K \in A$ y $L / H = K /H$. Probaremos que $L = K$. Sea $\ell \in L$. Entonces, $\ell H \in L / H$, i.e., existe $k \in K$ tal que $\ell H = k H$. Así, $\ell ^{-1} k \in H \subseteq K$. Como $k \in K$ y $\ell^{-1}k \in K$, tenemos $\ell = k \cdot (\ell^{-1}k)^{-1} \in K$, de donde $L \subseteq K$. La otra inclusión es igual. Luego $\Phi$ es inyectiva. Ahora, sea $C \in B$,  i.e, $C \leq G / H$. Sea $K = p ^{-1}(C) \leq G$ y $\Ker(p) \subseteq K$. Entonces, $H \subseteq K$. Así, $\Phi(K)=C$. Concluya que que $\Phi$ es una biyección. Probaremos ahora los puntos 1 y 2:
\begin{enumerate}[label=(\arabic*)]
\item Es claro que $L \leq K \iff L ^\ast \leq K^{\ast}$. Para probar que los indices coinciden, basta probar que existe una biyección entre las coclases de $K/L$ y $K^\ast/L^{\ast}$. Sea $f:K/L \to K^{\ast}/L^{\ast}$ tal que $kL \mapsto p(k)L^{\ast} = kH (L / H)$.
\begin{itemize}
\item \textbf{Bien-definición de $f$}: Si $k_{1}, k_{2} \in K$ son tales que $k_{1}L = k_{2}L$, entonces $k_{2}^{-1} k_{1} \in L$ y en consecuencia $k_{2}^{-1} k_{1} H \in L/H = L^{\ast}$, lo que implica que $f(k_{1}L) = p(k_{1}) L^\ast = p(k_{2}) L^{\ast} = f(k_{2}L)$, y por tanto, $f$ está bien definida.
\item \textbf{Inyectividad de $f$}: Sean $k_{1}, k_{2} \in K$ tales que $f(k_{1}L) = f(k_{2}L)$. Note que $f$ es inyectiva, pues
\[
\begin{aligned}
p(k_{1})L^{\ast} = p(k_{2}) L^{\ast} &\implies p(k_{2})^{-1} \ast p(k_{1}) \in L^{\ast} \implies p(k_{2}^{-1} k_{1}) \in L^\ast = L/H \\
k_{2}^{-1} k_{1} H \in L / H &\implies k_{2}^{-1} k_{1} \in L \implies k_{1}L = k_{2}L,
\end{aligned}
\]
\item \textbf{Sobreyectividad de $f$}: Esto se sigue directamente de la sobreyectividad de $p$.
Así, $[K:L] = [K^{\ast}:L^{\ast}]$.
\end{itemize}
\item El isomorfismo entre $K/L$ y $K^{\ast}/L^{\ast}$ se sigue del tercer teorema del isomorfismo.
\end{enumerate}
\end{proof}

% ======================================================================
\section{Acciones de Grupo y Grupos Simples}
% ======================================================================


\subsection{Definiciones y propiedades elementales}

\begin{namedstd}{Definición}{Acción de Grupos}
Sea $(G, \cdot)$ un grupo y sea $X$ un conjunto. Decimos que $G$ actúa sobre $X$ (o que $X$ es un $G$-conjunto) si existe una función $\alpha: G \times X \to X$ tal que
\begin{enumerate}[label=(\arabic*)]
\item $\forall g,h \in G  \quad\forall x \in X  \quad \Big( \alpha\big(g, \alpha(h,x)\big) = \alpha(g \cdot h, x)\Big)$,
\item $\forall x \in X  \quad\Big( \alpha(1_{G}, x) = x \Big)$.
\end{enumerate}

\end{namedstd}
\begin{namedstd}{Notación}{Acción de un elemento}
Si $G$ actúa sobre $X$, dado $g \in G$, denotamos $\alpha_{g}:X \to X$ tal que $x \mapsto \alpha(g,x)$. Si conocemos la acción específica, es común escribir $g \cdot x$.

\end{namedstd}
\begin{namedstd}{Nota}{Propiedades inmediatas}
\begin{enumerate}[label=(\arabic*)]
\item Para todos $g,h \in G$ y para todo $x \in X$, tenemos que $\alpha_{g}(\alpha_{h}(x)) = \alpha_{g \cdot h}(x)$.
\item $\alpha_{1_{G}} = \mathrm{id}_{X}$
\end{enumerate}

\end{namedstd}
\begin{namedstd}{Ejemplos}{Acciones usuales}
\begin{enumerate}[label=(\arabic*)]
\item Sea $n \in \mathbb{N}$. Considere $G = S_{n}$ y $X = \{ 1,\dots,n \}$. Defina $\alpha: S_{n} \times X \to X$ tal que $(\sigma, i) \mapsto \sigma(i)$. Note que si $\gamma,\sigma \in S_{n}$, dado $i \in X$,
\[
\alpha_{\sigma}(\alpha_{\gamma}(i)) = \alpha_{\sigma}(\gamma(i)) = (\sigma \circ \gamma)(i) = \alpha_{\sigma \circ \gamma} (i).
\]
\end{enumerate}
Además, $\alpha_{\mathrm{id}}(i) = \mathrm{id}(i) = i$. Por tanto, $\alpha$ es una acción de grupos y $S_{n}$ actúa sobre $X$.

\begin{enumerate}[label=(\arabic*)]
\item Sean $G = (\mathbb{Z}, +)$ y $X = \mathbb{R}$. Defina $\tau: G \times X \to X$ tal que $(n,x) \mapsto x+ n$. Note que dados $m,n \in \mathbb{Z}$ y $x \in \mathbb{R}$,
\[
\tau_{n}(\tau_{m(x)}) = \tau_{n}(x+m) = (x+m)+n = x + (m+n) = \tau_{m+n}(x).
\]
\end{enumerate}
Además, $\tau_{0}(x) = x+0 = x$. Por tanto, $\mathbb{Z}$ actúa sobre $\mathbb{R}$.

\begin{enumerate}[label=(\arabic*)]
\item Sean $G = \mathrm{GL}_{n}(\mathbb{R})$ y $X = \mathbb{R}^{n}$. Defina $\alpha: G \times X \to X$ tal que $(A,v) \mapsto A v$. Mostrar que $\alpha$ es una acción de grupos queda como ejercicio.
\end{enumerate}

\begin{enumerate}[label=(\arabic*)]
\item Dado un grupo $G$, defina $\mu:G \times G \to G$ tal que $(g,h) \mapsto g \cdot h$. Es fácil ver que $\mu$ es una acción de grupo. Así, $G$ actúa sobre sí mismo.
\end{enumerate}

\end{namedstd}
\begin{namedres}{Lema}{Acción es biyectiva sobre $X$}
Sea $G$ un grupo y $X$ un conjunto. Suponga que $\alpha:G \times X \to X$ es una acción de grupos. Entonces, si $g \in G$, $\alpha_{g}: X \to X$ tal que $x \mapsto \alpha(g,x)$ es una biyección, i.e., $\alpha_{g} \in S_{x}$.
\end{namedres}
\begin{proof}
Considere $\alpha_{g^{-1}}:X \to X$ tal que $x \mapsto \alpha(g^{-1}, x)$. Note que
\[
(\alpha_{g} \circ \alpha_{g^{-1}})(x) = \alpha_{g \circ g ^{-1}}(x) = \alpha_{1_{G}}(x) = x.
\]
De la misma forma, $(\alpha_{g^{-1}} \circ \alpha_{g}) = \mathrm{id}_{X}$, por lo que concluimos que la función es biyectiva.


\end{proof}
\begin{namedres}{Teorema}{Acciones y homomorfismos están en biyección}
Sea $G$ un grupo y $X$ un conjunto. Defina
\[
\begin{aligned}
\Omega:&= \{ \alpha: G \times X  \to X: \alpha \text{ es una acción de } G \text{ en }X\},\\
\Sigma:&= \{ \phi: G \to S_{X} : \phi \text{ es homomorfismo de grupos} \}.
\end{aligned}
\]
Entonces, $\Omega$ y $\Sigma$ están en biyección
\end{namedres}
\begin{proof}
Defina $\Phi: \Omega \to \Sigma$ de la siguiente manera. Dada una acción $\alpha: G \times X \to X$, defina $\phi:G \to S_{X}$ tal que $\phi(g) = \alpha_{g}$ donde $\alpha_{g}:X \to X$ es tal que $x \mapsto\alpha(g,x)$. Es fácil ver que $\phi$ es un homomorfismo y que $\alpha_{g} \in S_{X}$. Tome $\Phi(\alpha) = \phi$.
Ahora, defina $\Psi: \Sigma \to \Omega$ de la siguiente manera. Dado $\phi:G\to S_{X}$ homomorfismo, sea $\alpha:G \times X \to X$ tal que $(g,x) \mapsto \phi(g)(x) \in X$. Note que $\alpha$ es una acción, pues si $g,h \in G$y $x \in X$, tenemos que
\[
\begin{aligned}
\alpha(g \cdot h,x) = \phi(gh) (x) &= (\phi(g) \circ \phi(h))(x) \\
&= \phi(g) \big(\phi(h)(x)\big) \\
&=\phi(g) (\alpha(h, x)) \\
&=\alpha(g, \alpha(h,x)),
\end{aligned}
\]
y además, dado $x \in X$, $\alpha(1_{G}, x) = \phi(1_{G})(x) = \mathrm{id}_{X}(x) = x$. Tome $\Psi(\phi) = \alpha$. Verificar que $\Phi \circ \Psi = \mathrm{id}_{\Sigma}$ y que $\Psi \circ \Phi = \mathrm{id}_{\Omega}$ queda como ejercicio.

\end{proof}
\begin{namedres}{Teorema}{Cayley}
Dado $G$ un grupo, $G$ es isomorfo a un subgrupo de $S_{G}$.
\end{namedres}
\begin{proof}
Considere $\mu_{g}:G \to G$ tal que $h \mapsto g\cdot h$. Sea $\phi: G \to S_{G}$ tal que $g \mapsto \mu_{g}$. Ya vimos que $\phi$ es un homomorfismo. Veamos que $\phi$ es inyectivo:
\[
g \in \Ker(\phi) \iff \phi(g) = \mu_{g}= \mathrm{id}_{G} \iff g \cdot x = x  \quad \forall x \in G \iff g =1_{G},
\]
por lo que $\Ker(\phi) = \{ 1_{G} \}$. Ahora, usando el primer teorema del isomorfismo, obtenemos que
\[
G / \Ker(\phi) = G/\{ 1_{G} \} = G\cong \Imm(\phi) \leq  S_{G},
\]
de donde concluimos el resultado.

\end{proof}

\subsection{Órbitas y estabilizadores}

Sea $G$ un grupo y $X$ un conjunto. Sea $\alpha$ una acción de $G$ sobre $X$. Defina la relación $\sim$ en $X$ tal que $x \sim y \iff \exists g \in G \Big(\alpha_{g}(x) = y\Big)$. Es fácil ver que $\sim$ es una relación de equivalencia.

\begin{namedstd}{Definición}{Órbita}
Sea $G$ un grupo y $X$ un conjunto tal que $G$ actúa sobre $X$ por la acción $\alpha$. Dada la relación de equivalencia $\sim$ arriba definida, definimos la órbita de $x \in X$ como
\[
\mathcal{O}_{x} := \{ y \in X:y \sim x \} = \{ y \in X: \exists g \in G \big( \alpha_{g}(x) = y \big) \}.
\]
Como $\sim$ es una relación de equivalencia, las órbitas forman una partición de $X$, i.e., $X = \dot\bigcup_{x \in X} \mathcal{O}_{x}$.

\end{namedstd}
\begin{namedstd}{Definición}{Acción transitiva}
Sea $G$ un grupo y $X$ un conjunto tal que $G$ actúa sobre $X$. Si existe $x \in X$ tal que $\mathcal{O}_{x} =X$ decimos que la acción es transitiva. Equivalentemente, la acción es transitiva si para todos $x,y \in X$ existe $g \in G$ tal que $x = \alpha_{g}(y)$.

\end{namedstd}
\begin{namedstd}{Definición}{Estabilizador}
Sea $G$ un grupo y $X$ un conjunto tal que $G$ actúa sobre $X$ por la acción $\alpha$. Dado $x \in X$, definimos el estabilizador de $x$ por $G_{x}:= \{ g \in G:\alpha_{g}(x)=x \}$.

\end{namedstd}
\begin{namedres}{Lema}{Estabilizador es un subgrupo}
\end{namedres}
\begin{proof}
$G_{x} \neq \emptyset$ pues $1_{G} \in G_{x}$ por definición. Ahora, sean $g,h \in G_{x}$. Note que $\alpha_{h^{-1}}(x) = (\alpha_{h})^{-1}(x) = x$. Finalmente, tenemos que
\[
\alpha_{gh^{-1}}(x) = \alpha_{g}(\alpha_{h^{-1}}(x)) = x,
\]
por lo que $g \cdot h^{-1} \in G_{x}$, de donde concluimos que $G_{x} \leq G$.


\end{proof}
\begin{namedstd}{Ejemplos}{Órbitas y estabilizadores}
\begin{enumerate}[label=(\arabic*)]
\item Sea $n\geq2$. Considere la acción de $S_{n}$ sobre $X = \{ 1,\dots,n \}$ tal que $\alpha_{\sigma}(i) = \sigma(i)$. En este caso hay una única órbita. Sean $i,j \in X$ y considere $\tau = \begin{pmatrix}i & j\end{pmatrix}$. Note que $\alpha_{\tau}(i) = j$, por lo que $i \sim j$, i.e., todo par de elementos de $X$ se relacionan. Por tanto la acción es transitiva. Ahora, note que si $1 \leq i\leq n$, $G_{i} = \{ \sigma \in S_{n}: \sigma(i) = i \} \cong S_{n-1}$.
\end{enumerate}

\begin{enumerate}[label=(\arabic*)]
\item Considere $G = (\mathbb{Z}, + )$ y $X=\mathbb{R}$ bajo la acción $\tau(n,x) = x+n$. Note que si $x,y \in \mathbb{R}$,
\[
x \sim y \iff \exists n \in \mathbb{Z} \Big(\tau_{n}(x) = y\Big) \iff \exists n \in \mathbb{Z} \Big(y=x+n\Big) \iff x-y \in \mathbb{Z},
\]
\end{enumerate}
de donde tenemos que $x$ y $y$ tienen la misma parte decimal. Hay una cantidad no numerable de órdenes. Es evidente ver que para todo $x \in \mathbb{R}$, $G_{x} = \{ 0 \}$.

\begin{enumerate}[label=(\arabic*)]
\item Considere $G = \mathbb{S}^{1} =\{ e^{i\theta}:\theta \in [0, 2\pi) \}$ y $X = \mathbb{C}$. Defina la acción $\rho_{\theta}(z) = z e^{i \theta}$ (ver que esto es una acción es un ejercicio). Sea $z \in \mathbb{C}$. Entonces,
\[
\mathcal{O}_{z} = \{ \rho_{\theta}(z): \theta \in [0, 2\pi) \} = \{ ze^{i \theta}: \theta \in [0,2\pi) \} = \{ w \in \mathbb{C}: \lvert w \rvert  = \lvert z \rvert  \}.
\]
\end{enumerate}
Ahora, veamos qué pasa con el estabilizador. Si $z = 0$, tenemos que
\[
G_{0} = \{ \theta \in \mathbb{S}^{1}: \rho_{\theta}(0) = 0 \} = \{ g \in \mathbb{S}^{1}: 0 e^{i \theta} = 0 \} = \mathbb{S}^{1}.
\]
Por otro lado, si $z \neq 0$,
\[
\rho_{\theta}(z) = z \iff z e^{i \theta} = z \iff \theta = 0 \iff e^{i \theta } = 1,
\]
lo que implica que $G_{Z} = \{ 1 \}$.

\begin{enumerate}[label=(\arabic*)]
\item Sea $G$ un grupo y $H \leq G$. Considere $G / H = \{ gH : g \in G\}$. Defina una acción de $G$ en $G / H$ tal que $\mu:G \times G / H\to G /H$ tal que $\mu_{g}(hH) = (g \cdot h) H$. Queda como ejercicio ver que $\mu$ es una acción transitiva. Ahora, note que
\[
g \in G_{xH} \iff (g \cdot x) H = xH \iff x ^{-1} \cdot g \cdot x \in H \iff g \in x H x ^{-1},
\]
\end{enumerate}
de modo que $G_{xH} = xH x ^{-1}$.

\begin{enumerate}[label=(\arabic*)]
\item Sea $G = \mathrm{GL}(n,\mathbb{R})$ y $X = \mathbb{R}^n$. Defina la acción $\alpha_{A}(v)=Av$. Note que
\[
\mathcal{O}_{0} = \{ v \in \mathbb{R}: \exists A \in G (Av = w) \} = \{ 0 \}.
\]
\end{enumerate}
Por tanto, esta acción no es transitiva.

\begin{enumerate}[label=(\arabic*)]
\item Sea $G$ un grupo. Defina $\mu:G\to G$ tal que $\mu_{g}(h) = gh$. Sean $h,g \in G$. Note que $\mu_{gh^{-1}}(h) = (gh^{-1})h=g$, de modo que $h \sim g$. Por tanto, la acción es transitiva.
\end{enumerate}

\begin{enumerate}[label=(\arabic*)]
\item Sea $G$ un grupo y considere la acción de $G$ en $G$ como la conjugación $\mu_{g}(h) = ghg^{-1}$. Sea $x \in G$. Tenemos que
\[
x^G := O_{x} = \{ y \in G: \exists g \in G (g xg^{-1} = y) \} = \{ gxg^{-1}:g \in G \}.
\]
\end{enumerate}
Por otro lado, $G_{x}:= \{ g \in G: gxg^{-1} = x \} = C_{G}(x)$ (el centralizador).

\end{namedstd}
\begin{namedres}{Teorema}{Órbita y cociente del estabilizador}
Sea $G$ un grupo y $X$ un conjunto tal que $G$ actúa sobre $X$. Sea $x \in X$. Entonces $\lvert \mathcal{O}_{x} \rvert = [G:G_{x}]$.

\end{namedres}
\begin{proof}
Vamos a construir una biyección $\Phi:\mathcal{O}_{x} \to G / G_{x}$. Sea $y \in \mathcal{O}_{x}$. Entonces existe $g \in G$ tal que $\alpha_{g}(x) = y$. Defina $\Phi(y) = g \cdot G_{x}$. $\Phi$ está bien definida, pues si $g,h$ son tales que $\alpha_{g}(x) = \alpha_{h}(x)$, entonces
\[
\alpha_{h^{-1}g}(x) = \alpha_{h^{-1}}(\alpha_{g}(x)) = \alpha_{h^{-1}}(\alpha_{h}(x)) = \alpha_{1_{G}}(x) = x,
\]
de modo que $h^{-1}g \in G_{x}$ y por tanto $g \cdot G_{x} = h \cdot G_{x}$. $\Phi$ es inyectiva pues si $y,z \in \mathcal{O}_{x}$ tal que $\Phi(y) = \Phi(z)$, existen $g,h \in G$ tales que $\alpha_{g}(x) = y$ y $\alpha_{h}(x) = z$. Así, tenemos que
\[
g G_{x} = h G_{x} \implies h^{-1} g \in G_{x} \implies \alpha_{h^{-1}g}(x)=x.
\]
Así, $\alpha_{h^{-1}}(\alpha_{g}(x)) = \alpha_{h}^{-1}(\alpha_{g}(x)) = x \implies \alpha_{g}(x)=\alpha_{h}(x)$. Finalmente $\Phi$ es sobreyectiva, pues si $g \cdot G_{x} \in G / G_{x}$, tome $y = \alpha_{g}(x)$. Note que $\Phi(y) = g \cdot G_{x}$. Así, $\Phi$ es una biyección y por tanto los conjuntos dados poseen la misma cardinalidad.

\end{proof}
\begin{namedres}{Corolario}{La órbita divide al orden del grupo}
Sea $G$ un grupo finito y $X$ un conjunto tal que $G$ actúa sobre $X$. Entonces para todo $x \in X$, $\lvert \mathcal{O}_{x} \rvert \Big\lvert \lvert G \rvert$.
\end{namedres}
\begin{proof}
Tenemos que $\lvert \mathcal{O}_{x} \rvert = [G:G_{x}] =\frac{ \lvert G \rvert}{\lvert G_{x} \rvert}$, de donde concluimos el resultado

\end{proof}
\begin{namedres}{Corolario}{Tamaño de la clase de conjugación}
Sea $G$ un grupo finito y $x \in G$, Entonces $\lvert x^{G} \rvert = [G:C_{G}(x)]$.

\end{namedres}
\begin{namedres}{Teorema}{Cauchy}
Sea $G$ un grupo finito y $p$ un primo tal que $p \Big\lvert \lvert G \rvert$. Entonces existe $g \in G$ tal que $\lvert g \rvert = p$.

\end{namedres}
\begin{proof}[\textbf{Prueba 1, para grupos abelianos.}]
Suponga que $G$ es abeliano. Procedemos por inducción fuerte sobre $\lvert G \rvert=n$. El caso base ($n=1$) es trivialmente cierto por vacuidad. Suponga como hipótesis inductiva que el resultado es cierto para todo grupo $H$ tal que $\lvert H \rvert < n$. Sea $a \in G$ con $a \neq 1_{G}$ y $k = \lvert  a \rvert$. Considere los siguientes casos:
\textbf{Caso 1:} Suponga que $p \mid k$. Entonces, existe $\ell \in \mathbb{N}$ tal que $p \ell = k$. Luego $\lvert a^{\ell} \rvert = p$.
\textbf{Caso 2:} Suponga que $p \not\mid k$. Sea $H = \langle a \rangle$. Como $G$ es abeliano, $H \triangleleft G$, por lo que $G/H$ es un grupo y $\lvert H \rvert = k$. Note que
\[
\lvert G / H \rvert = \frac{\lvert G \rvert }{\lvert H \rvert } = \frac{n}{k} < n.
\]
Como $p \mid n$ y $p \not\mid k$, $p \mid \frac{n}{k}$. Por la hipótesis inductiva, existe $sH \in G /H$ tal que $\lvert sH \rvert = p$. Sea $m = \lvert s \rvert$. Note que $(sH)^{m} = s^{m}H = H$, por lo que $p \mid m$. Finalmente, aplique el caso 1 sobre $s$. Concluimos el resultado

\end{proof}
\begin{proof}[\textbf{Prueba 2, general.}]
Sea $x \in G$. Recuerde que $x^{G} = \{ g x g^{-1}: g \in G \}$ y que $\lvert x^{G} \rvert = [G:C_{G}(x)]$. Procedemos por inducción fuerte sobre $\lvert G \rvert=n$.  Considere los siguientes casos:
\textbf{Caso 1:} Si $x \in Z(G)$, entonces $x^{G}=\{ x \} \implies G = C_{G}(x)$.
\textbf{Caso 2:} Si $x \not\in Z(G)$, entonces $\lvert x^{G} \rvert > 1$.  Entonces $C_{G}(x) < G \implies \lvert C_{G}(x) \rvert < \lvert G \rvert$. Si $p \Big\lvert \lvert C_{G}(x) \rvert$, por la hipótesis inductiva, existe $g \in C_{G}(x) \subseteq G$ tal que $\lvert g \rvert=p$, de donde concluimos el resultado. Si $p \not\Big\lvert \lvert C_{G}(x) \rvert$, tenemos que
\[
[G:C_{G}(x)] = \frac{\lvert G \rvert }{\lvert C_{G}(x) \rvert } \implies [G:C_{G}(x)] \cdot \lvert C_{G}(x) \rvert = \lvert G \rvert \implies p \mid [G:C_{G}(x)].
\]
Sean $x_{1},\dots,x_{k}$ las clases de conjugación diferentes que se obtienen mediante la relación de ser conjugación y que no estén en $Z(G)$. Entonces,
\[
G = Z(G) \dot{\cup} \bigcup_{i=1}^{k} x_{i}^{G} \implies \lvert G \rvert = \lvert Z(G) \rvert + \sum_{i=1}^{k} \lvert x_{i}^{G} \rvert \implies \lvert G \rvert - \sum_{i=1}^{k} \lvert x_{i}^{G} \rvert = \lvert Z(G) \rvert, 
\]
de donde concluimos que $p \Big\lvert \lvert Z(G) \rvert$. Así,  $Z(G)$ es subgrupo propio de $G$, y concluimos el resultado al aplicar la hipótesis inductiva sobre $Z(G)$.

\end{proof}
\begin{namedstd}{Nota}{La hipótesis de primalidad es necesaria}
El teorema es falso si $p$ no es primo. Note que $8 \mid 120 \implies 8 \Big\lvert \lvert S_{5} \rvert$. Pero hemos visto que no hay subgrupos de orden 8 en $S_{5}$.


\end{namedstd}
\begin{namedstd}{Definición}{$p$-grupo}
Sea $p$ primo. Un grupo $G$ es un $p$-grupo si existe $k \in \mathbb{N}$ tal que $\lvert G \rvert=p^{k}$.

\end{namedstd}
\begin{namedstd}{Ejercicio}{Centro de un $p$-grupo}
Si $G$ es un $p$-grupo, $Z(G) \neq \{ 1_{G} \}$.

\end{namedstd}
\begin{namedres}{Corolario}{Grupos de orden $p^{2}$}
Si $p$ es primo y $\lvert G \rvert = p^{2}$, entonces $G$ es abeliano.
\end{namedres}
\begin{proof}
Suponga que $G$ no es abeliano. Entonces $Z(G)$ es un subgrupo propio de $G$. Por Lagrange, $\lvert Z(G) \rvert \in \{ 1,p \}$. Por el ejercicio anterior, $\lvert Z(G) \rvert = p$. Sabemos que $Z(G) \triangleleft G$. Así, $\lvert G / Z(G) \rvert = \frac{\lvert G \rvert}{\lvert Z(G) \rvert} = p$, de donde concluimos que $G / Z(G)$ es cíclico y por tanto $G$ es abeliano.


\end{proof}
\begin{namedres}{Teorema}{Subgrupos de un grupo abeliano finito}
Sea $G$ un grupo abeliano finito. Entonces $G$ tiene un subgrupo de orden $d$ para todo $d$ tal que $d \Big\lvert \lvert G \rvert$.
\end{namedres}
\begin{proof}
Sea $n = \lvert G \rvert$. Proceda por inducción fuerte sobre $d$, con $d \mid n$. El caso base es trivial, pues si $H = \{ 1_{G} \}\implies \lvert H \rvert=1$. Sea $d>1$. Suponga como hipótesis inductiva que para todo $e < d$ y para todo $\tilde{G}$ grupo con $e \Big\lvert \lvert \tilde{G} \rvert$ existe $\tilde{H} < \tilde{G}$ tal que $\lvert \tilde{H} \rvert = e$. Como $d \mid n$, existe $k \in \mathbb{N}$ tal que $kd=n$. Como $d>1$, existe $p$ primo tal que $p \mid d$, i.e., existe $\ell \in \mathbb{N}$ tal que $\ell p = d$. Por Teorema de Cauchy, existe $g \in G$ tal que $\lvert g \rvert=p$. Sea $H = \langle g \rangle$. Note que $H$ es cíclico, y por tanto abeliano y normal. Luego, $G / H$ es grupo. Así,
\[
\lvert G / H \rvert = \frac{\lvert G \rvert }{\lvert H \rvert } = \frac{n}{p} = \frac{kd}{p} = \frac{k \ell p}{p} = k \ell.
\]
Así, $\ell \Big\lvert \lvert G / H \rvert$ y $\lvert G/H \rvert < n$. Como $\ell < d$, por la hipótesis inductiva, existe $S^{\ast} < G / H$ tal que $\lvert S^{\ast} \rvert = \ell$. Por teorema de correspondencia, existe $S < G$ tal que $H \leq S \leq G$ y $S^{\ast} = S/H$. Concluimos notando que
\[
\ell = \lvert S^{\ast} \rvert = \frac{\lvert S \rvert}{\lvert H \rvert } = \frac{\lvert S \rvert}{p} \implies \lvert S \rvert = \ell p =d, 
\]
de donde se sigue el resultado.

\end{proof}
\begin{namedres}{Teorema}{Subgrupos de un $p$-grupo}
Sea $G$ un $p-grupo$ tal que $\lvert G \rvert = p^{e}$. Entonces para todo $k \leq e$, existe $H \leq G$ tal que $\lvert H \rvert = p^{k}$.

\end{namedres}
\begin{proof}
Procedemos por inducción sobre $e$. Para $e=0$, tenemos que $G=\{ 1_{G} \}.$ Tome $H = G$. Sea $e > 0$ y sea $k \leq e$.  Suponga como hipótesis inductiva que el resultado es cierto para $f<e$. Como $\lvert G \rvert = p^{e}$ y $e>0$, entonces $Z(G)$ es no trivial. Considere los siguientes casos.
\begin{itemize}
\item Si $G = Z(G)$, entonces $G$ es abeliano. Por el teorema anterior, concluimos el resultado.
\item Ahora, si $G \neq Z(G)$, existe $c < e$ tal que $\lvert Z(G) \rvert = p^{c}$ con $c<e$. Si $k \leq c$, entonces por hipótesis inductiva sobre $Z(G)$, existe $H \triangleleft Z(G)$ tal que $\lvert H \rvert=p^{k}$. Por otro lado, si $k>c$, sabemos que $G / Z(G)$ es un grupo (pues $Z(G) \triangleleft G$). Ahora, note que
\[
\lvert G / Z(G) \rvert = \frac{\lvert G \rvert }{\lvert Z(G) \rvert } = p^{e-c}.
\]
\end{itemize}
Aplicando hipótesis inductiva sobre $G / Z(G)$, existe $S^{\ast} \leq G/Z(G)$ tal que $\lvert S^{\ast} \rvert = p^{k-c}$. Por teorema de correspondencia, existe $S \leq G$ tal que $Z(G) \leq S \leq G$ y $S^{\ast} = S / Z(G)$. Así,
\[
\lvert S^{\ast} \rvert \cdot \lvert Z(G) \rvert = \lvert S \rvert \implies p^{k-c} \cdot p^{c} = \lvert S \rvert,
\]
de donde concluimos el resultado.

\end{proof}

\subsection{Grupos Simples}


\begin{namedstd}{Definición}{Grupo simple}
Sea $G$ un grupo no trivial. Decimos que $G$ es simple si no tiene subgrupos normales propios, i.e., si $H \triangleleft G$, entonces $H = \{ 1_{G} \}$ o $H=G$.

\end{namedstd}
\begin{namedres}{Teorema}{Abelianos simples son $\mathbb{Z}_{p}$}
Sea $G$ un grupo abeliano. Entonces $G$ es simple si y solo si $G=\mathbb{Z}_{p}$ con $p$ primo.

\end{namedres}
\begin{proof}
($\impliedby$): Esta dirección es sencilla, pues basta aplicar Lagrange para notar que $G$ no tiene subgrupos propios.
($\implies$): Como $G$ es abeliano, todo subgrupo es normal. Además, como $G$ es simple, $G \neq \{ 1_{G} \}$. Entonces, existe $x \in G$ tal que $(x \neq 1_{G})$. Tenemos que
\[
\langle x \rangle \triangleleft  G \implies \langle x \rangle = G \implies G \text{ es cíclico} \implies G \cong \mathbb{Z} \quad \lor \quad G \cong   \mathbb{Z}_{n}, \text{ para algún } n \in \mathbb{N}.
\]
Como $G$ es simple, $G \not\cong \mathbb{Z}$, por lo que $G \cong \mathbb{Z}_{n}$ para algún $n \in  \mathbb{N}$. Además, como $G$ es simple, $n$ es primo, de donde concluimos el resultado.

\end{proof}
\begin{namedstd}{Recordatorio}{Signo y grupo alternante}
Si $\sigma \in S_{n}$ tal que $\sigma = \beta_{1} \cdots \beta_{t}$ es su descomposición en ciclos disjuntos de $\sigma$, entonces, $\sgn( \sigma ) := (-1)^{n-t}$. Si $\sgn( \sigma ) = 1$ decimos que $\sigma$ es par y $A_{n} = \{ \sigma \in S_{n}, \sigma \text{ es par} \}.$

\end{namedstd}
\begin{namedstd}{Ejercicios}{Sobre $3$-ciclos}
\begin{enumerate}[label=(\arabic*)]
\item $A_{n}$ es generado por los 3-ciclos si $n\geq {5}$.
\item En $S_{5}$, todos los 3-ciclos son conjugados.
\item En $A_{5}$, todos los 3-ciclos son conjugados.
\end{enumerate}
\end{namedstd}
\begin{namedres}{Teorema}{$A_{5}$ es simple}
$A_{5}$ es simple

\end{namedres}
\begin{proof}
Sea $H \neq \{ \mathrm{id}_{S_{5}} \}$ tal que $H \triangleleft A_{5}$. Por los ejercicios anteriores, basta mostrar que $H$ contiene un 3-ciclo, pues si $\sigma \in H$ es un 3-ciclo, $\sigma^{A_{5}} \subseteq H$ pues $H \triangleleft A_{5}$, entonces $\langle \sigma^{A_{5}}  \rangle = \{ \tau: \tau \text{ es 3-ciclo}  \} \subseteq H$ y por tanto $H = A_{5}$.

¿Cuántas configuraciones posibles tengo para la descomposición en ciclos disjuntos?
\begin{enumerate}[label=(\arabic*)]
\item $\sigma$ se descompone en 5 1-ciclos, y por tanto $\sigma \in A_{5}$.
\item $\sigma$ se descompone en un 2-ciclo y 3 1-ciclos, y por tanto $\sigma \not\in A_{5}$.
\item $\sigma$ se descompone en un 3-ciclo y 2 1-ciclos, y por tanto $\sigma \in A_{5}$.
\item $\sigma$ se descompone en un 4-ciclo y 1 1-ciclo, y por tanto $\sigma \not\in A_{5}$.
\item $\sigma$ se descompone en un 5-ciclo, y por tanto $\sigma \in A_{5}$.
\item $\sigma$ se descompone en dos 2-ciclos y un 1-ciclo, y por tanto $\sigma \in A_{5}$.
\item $\sigma$ se descompone en un 3-ciclo y un 2-ciclo, y por tanto $\sigma \not\in A_{5}$.
\end{enumerate}
Luego, si $\sigma \in H \setminus \{ 1_{G} \}$, su factorización es del tipo 3, 5 o 6. Si $\sigma$ es del tipo 3, entonces posee un 3-ciclo y por tanto se cumple el resultado.
SI $\sigma$ es del tipo 5, suponga sin pérdida de generalidad que $\sigma = \begin{pmatrix}1 & 2 & 3 & 4 & 5\end{pmatrix}$. Sea $\alpha = \begin{pmatrix}1 & 2 & 3\end{pmatrix} \in A_{5}$. Luego,
\[
w = \begin{pmatrix}
2 & 3 & 1 & 4 & 5
\end{pmatrix} =\alpha \sigma \alpha^{-1} \in H, \text{ pues }H \triangleleft A_{5} \implies w \sigma ^{-1} = \begin{pmatrix}
1 & 2 & 4
\end{pmatrix} \in H.
\]
Si $\sigma$ es del tipo 6, suponga sin pérdida de generalidad que $\sigma = \begin{pmatrix}1 & 2\end{pmatrix}\begin{pmatrix}3 & 4\end{pmatrix}$. Sea $\beta = \begin{pmatrix}3 & 4 & 5\end{pmatrix} \in A_{5}$. Note que $\delta :=\beta \sigma \beta ^{-1} = \begin{pmatrix}1 & 2\end{pmatrix}\begin{pmatrix}4 & 5\end{pmatrix} \in H$. Finalmente, $\sigma \cdot \delta = \begin{pmatrix}3 & 4 & 5\end{pmatrix} \in H$.
Concluya que $A_{5}$ es simple.

\end{proof}
\begin{namedres}{Teorema}{$A_{6}$ es simple}
$A_{6}$ es simple

\end{namedres}
\begin{proof}
Sea $H\neq \{ \mathrm{id}_{S_{6}} \}$ tal que $H \triangleleft A_{6}$. Suponga por contradicción que si $\sigma \in H \setminus \{ \mathrm{id}_{S_{6}} \}$, entonces $\sigma$ mueve todo elemento de $\{ 1,2,\dots,6 \}$.  ¿Cuántas configuraciones posibles tengo para la descomposición en ciclos disjuntos?
\begin{enumerate}[label=(\arabic*)]
\item $\sigma$ se descompone en tres 2-ciclos, y por tanto $\sigma \not\in A_{6}$.
\item $\sigma$ se descompone en un 2-ciclo y un 4-ciclo, y por tanto $\sigma \in A_{6}$.
\item $\sigma$ se descompone en dos 3-ciclos, y por tanto $\sigma \in A_{6}$.
\end{enumerate}
Si $\sigma$ es del tipo 2, considere $\sigma = \begin{pmatrix}1 & 2\end{pmatrix}\begin{pmatrix}3 & 4 & 5 & 6\end{pmatrix}$. Como $H$ es subgrupo, si $\sigma \in H$, $\sigma^{2} \in H$. Pero $\sigma^{2}(1) = 1$, una contradicción.
Por otro lado, si $\sigma$ es del tipo 3, considere $\sigma=\begin{pmatrix}1 & 2 & 3\end{pmatrix}\begin{pmatrix}4 & 5 & 6\end{pmatrix}$. Luego, tenemos que $\sigma ^{-1} = \begin{pmatrix}3 & 2 & 1\end{pmatrix}\begin{pmatrix}6 & 5 & 4\end{pmatrix}$. Sea $\gamma = \begin{pmatrix}2 & 3 & 4\end{pmatrix}$. Como $H$ es subgrupo, tenemos que $\sigma\gamma \sigma ^{-1} \gamma ^{-1} = \begin{pmatrix}1 & 5 & 3 & 2 & 4\end{pmatrix} \in H$, una contradicción.
Concluya que existen $\sigma \in H$ y $i \in \{ 1,\dots,6 \}$ tales que $\sigma(i) = i$. Sea $F = \{ \sigma \in A_{6}: \sigma(i)=i \}$. Note que, $F \neq \{ 1_{S_{6}} \}$ y $F \cap H \neq \{ 1_{S_{6}} \}$. Es fácil ver que $F \cong A_{5}$. Finalmente, como $F \cap H \triangleleft F \cong A_{5}$ y $A_{5}$ es simple, tenemos que $F \cap H = F$, de donde $F \subseteq H$. Concluya que $A_{6}$ es simple porque contiene un 3-ciclo.
\end{proof}
\begin{namedres}{Teorema}{$A_{n}$ es simple para $n \geq 5$}
$A_{n}$ es simple para $n \geq5$

\end{namedres}
\begin{proof}
Considere el caso para $n>6$. Sea $H\neq \{ \mathrm{id}_{S_{n}} \}$ tal que $H \triangleleft A_{n}$ Sea $\sigma \in H \setminus \{ \mathrm{id}_{Sn} \}$. Entonces, existe $i \in \{ 1,\dots,n \}$ tal que $\sigma(i) = j \neq i$. Sea $\alpha$ un 3-ciclo tal que $\alpha(i)=i$ y $\alpha(j)\neq j$. Note que
\[
\alpha(\sigma(i)) = \alpha(j) \neq  j  \quad \land  \quad\sigma(\alpha(i)) = \sigma(i) =j \implies \alpha \sigma \neq \sigma\alpha \implies \alpha\sigma\alpha ^{-1}\sigma ^{-1} \in H \setminus \{1_{S+n}\}.
\]
Sabemos que la conjugación de permutaciones no me cambia la estructura de ciclos. Luego $\sigma \alpha ^{-1} \sigma ^{-1}$ es un 3-ciclo y $\gamma = \alpha\sigma\alpha ^{-1}\sigma ^{-1} \in H$ es un producto de dos 3-ciclos y por tanto mueve a lo sumo 6 elementos. Sean $i_{1},\dots i_{k}$, con $k \leq 6$, los elementos que mueve $\gamma$. Ahora, sea $F:=\{ \sigma \in A_{n}: \sigma(\ell) = \ell  \quad \forall \ell \not\in \{ i_{1},\dots,i_{k} \}\}$. Es fácil ver que $F \cong S_{6}$. y $F \cap H \neq \{ 1_{S_{n}} \}$. Así, $F \cap H \triangleleft F \cong A_{6}$, por lo que $F \cap H = F$ y por tanto $F \subseteq H$. En particiular, $\begin{pmatrix}i_{1} & i_{2} & i_{3}\end{pmatrix} \in F \subseteq H$, de donde concluimos que es simple porque contiene un $3-$ciclo..

\end{proof}

\subsection{Teoremas de Sylow}


\begin{namedstd}{Definición}{Subgrupo de Sylow}
Sea $p$ un primo y $G$ un grupo. Decimos que $H\leq G$ es un $p-$subgrupo de Sylow si $\lvert H \rvert = p^{e}$ y $p^{e}$ es la mayor potencia de $p$ que divide a $\lvert G \rvert$.

\end{namedstd}
\begin{namedstd}{Ejemplo}{Sylow en $S_{4}$}
Considere $S_{4}$. Note que $\lvert S_{4} \rvert=4! = 2^{3} \cdot 3$. Un $2-$subgrupo de Sylow tiene tamaño $2^{3}$. Un $3-$subgrupo de Sylow tiene tamaño $3$. Un $5$-subgrupo de Sylow es $\{ e \}$.

\end{namedstd}
\begin{namedres}{Teorema}{Primer teorema de Sylow}
Si $p$ es primo y $G$ es finito, existe un $p-$subgrupo de Sylow.

\end{namedres}
\begin{proof}
Sea $\lvert G \rvert = p^{e}k$ tal que $p \not\mid k$. Sea $\Omega:=\{ T \subseteq G: \lvert T \rvert=p^{e} \}$. Considere la acción $\alpha$ de $G$ en $\Omega$ tal que $\alpha _g(T) = gT$. (queda como ejercicio ver que es una acción). Tenemos que $\lvert \Omega \rvert = \binom{p^{e}k}{p^{e}} \equiv k \mod p$. Como $p$ no divide a $k$, tenemos que $p\ \not\mid \Omega$. Sabemos que
\[
\Omega = \dot{\bigcup}_{i \in I} \mathcal{O}_{i},
\]
donde $\mathcal{O}_{i}$ son las órbitas distintas. Luego $\lvert \Omega \rvert = \sum_{i \in I} \lvert \mathcal{O}_{i} \rvert$. Por lo tanto, existe una órbita $\mathcal{O}_{i}$ tal que $p \not\mid \lvert \mathcal{O}_{i} \rvert$. Luego, existe $H \in \Omega$ tal que $p \not\mid \mathcal{O}_{H}$ (tomando $H$ como un representante de la clase de $\mathcal{O}_{i}$) . Tenemos que $\lvert \mathcal{O}_{H} \rvert = [G:G_{H}] = \frac{\lvert G \rvert}{\lvert G_{H} \rvert} = \frac{p^{e}k}{\lvert G_{H}. \rvert}$ De aquí, tenemos que $\lvert \mathcal{O}_{H} \rvert \mid p^{e}k \implies \lvert \mathcal{O}_{H} \rvert \mid k$. Sea $A = \bigcup_{T \in \mathcal{O}_{H}} T = \bigcup_{g \in G} gH$. Probaremos que $G=A$. Es claro que $A \subseteq G$.
Vamos a mostrar primero que $A$ es estable por traslaciones izquierdas. Sea $g \in G$, $T \in \mathcal{O}_{H}$. Hay que mostrar que $gA \subseteq A$. Dado $x \in A$, existe $T \in \mathcal{O}_{H}$ tal que $x \in T$ y por tanto, $gx \in gT \subseteq A$, de donde concluimos el resultado. Finalmente, tome $h_{0} \in H$ (existe pues $H \neq \emptyset$); entonces $1_{G} \in h_{0}^{-1}H \in \mathcal{O}_{H}$, de modo que $1_{G} \in A$. Dado cualquier $g \in G$, la estabilidad por traslaciones da $g = g \cdot 1_{G} \in gA \subseteq A$. Concluya que $G \subseteq A$ y por tanto $G=A$. Así, tenemos que
\[
p^{e}k = \lvert G \rvert = \lvert A \rvert =\left\lvert  \bigcup_{T \in \mathcal{O}_{H}}T  \right\rvert \leq \sum_{T \in \mathcal{O}_{H}}  \lvert T \rvert = \lvert \mathcal{O}_{H} \rvert \cdot \lvert T \rvert = \lvert \mathcal{O}_{H} \rvert \cdot p^{e},
\]
y por tanto, $k\leq \lvert \mathcal{O}_{H} \rvert$. Pero $\lvert \mathcal{O}_{H} \rvert \mid k$, de donde $\lvert \mathcal{O}_{H} \rvert = k$. Finalmente, como
\[
\lvert \mathcal{O}_{H} \rvert \cdot \lvert G_{H} \rvert = p^{e} \cdot k \implies \lvert G_{H} \rvert = \frac{p^{e}k}{k} = p^{e}.
\]
Así, $\lvert G_{H} \rvert = p^{e}$ y es un grupo.
\end{proof}
\begin{namedres}{Teorema}{Segundo teorema de Sylow}
Sea $G$ un grupo finito y $p$ primo. Si $P$ y $Q$ son $p-$grupos de Sylow, entonces son conjugados, i.e., existe $g \in G$ tal que $gPg^{-1} = Q$.
\end{namedres}
\begin{proof}
Ejercicio.
\textbf{Sugerencia:} Considere la acción $\alpha:P \times G / Q \to G / Q$ tal que $\alpha_{g}(hQ) = ghQ$.
\end{proof}
\begin{namedres}{Teorema}{Tercer teorema de Sylow}
Sea $G$ finito y $p$ primo y $n_{p}$ el número de $p-$grupos de Sylow. Entonces $p \mid n_{p}-1$. 
\end{namedres}
\begin{proof}
\textbf{Sugerencia:} Definir $\Lambda_{p}: \{ H \subseteq G: H \text{ es un }p-\text{subgrupo de Sylow} \}$. Considerar la acción $\alpha: G \times\Lambda_{p} \to \Lambda_{p}$ tal que $\alpha_{g}(P) = gPg^{-1}$.
\end{proof}



\appendix
% =============================================================================
\section{Anexo: caracterizaciones equivalentes}
% =============================================================================

A lo largo del curso, varios conceptos clave admitieron \emph{múltiples caracterizaciones equivalentes} que conviene tener juntas a la vista. Este anexo recopila esos paquetes de equivalencias en cajas autocontenidas, una por concepto. Cada caja indica además dónde se prueba cada caracterización dentro del documento.

\anexogroup{Subgrupos y generación}

\begin{equivbox}{blue}{Caracterizaciones de subgrupo $H \leq G$}
Sea $(G,\ast)$ un grupo y $\emptyset \neq H \subseteq G$. Las siguientes son equivalentes:
\begin{enumerate}[label=(\arabic*),leftmargin=2em]
\item \textbf{(definición)}\; $(H, \ast\mid_{H})$ es un grupo. \hfill\textit{§1.5.}
\item \textbf{(criterio $xy^{-1}$)}\; Para todos $x, y \in H$: $x \ast y^{-1} \in H$. \hfill\textit{Teorema 1.5.2.}
\item \textbf{(tres condiciones)}\; $H$ es cerrado bajo $\ast$, $1_{G} \in H$, y $h \in H \implies h^{-1} \in H$. \hfill\textit{Nota 1.5.1.}
\end{enumerate}
\end{equivbox}

\begin{equivbox}{green!60!black}{Caracterizaciones del subgrupo generado $\langle S \rangle$ y de los grupos cíclicos}
Sea $S \subseteq G$, $S \neq \emptyset$. Las siguientes describen el mismo objeto $\langle S \rangle$:
\begin{enumerate}[label=(\arabic*),leftmargin=2em]
\item \textbf{(intersección)}\; $\langle S \rangle = \bigcap \{ H : H \leq G \ \land\ S \subseteq H \}$. \hfill\textit{Definición 1.5.4.}
\item \textbf{(minimalidad)}\; $\langle S \rangle$ es el subgrupo más pequeño que contiene a $S$. \hfill\textit{Teorema 1.5.5.}
\item \textbf{(palabras)}\; $\langle S \rangle = \{ x_{1}^{\alpha_{1}} \ast \dots \ast x_{n}^{\alpha_{n}} : n \in \mathbb{N}^{\ast},\ x_{i} \in S,\ \alpha_{i} \in \{-1,1\} \}$. \hfill\textit{Teorema 1.5.6.}
\end{enumerate}
Para grupos cíclicos ($G = \langle g \rangle$):
\begin{enumerate}[label=(\arabic*),leftmargin=2em,resume]
\item \textbf{(clasificación)}\; $\lvert g \rvert = \infty \implies G \cong \Z$; \quad $\lvert g \rvert = n \implies G \cong \Z_{n}$. \hfill\textit{Teorema 2.1.18.}
\item \textbf{(herencia)}\; Todo subgrupo de un grupo cíclico es cíclico. \hfill\textit{Teorema 1.7.7.}
\end{enumerate}
\end{equivbox}

\anexogroup{Órdenes y generadores}

\begin{equivbox}{cyan!70!black}{Caracterizaciones del orden de un elemento}
Sea $g \in G$ con $\lvert g \rvert = n < \infty$ y $k \in \mathbb{Z}$. Entonces:
\begin{enumerate}[label=(\arabic*),leftmargin=2em]
\item \textbf{(definición)}\; $n$ es el menor natural positivo con $g^{n} = 1_{G}$. \hfill\textit{Definición 1.6.1.}
\item \textbf{(divisibilidad)}\; $g^{k} = 1_{G} \iff n \mid k$. \hfill\textit{Teorema 1.6.3.}
\item \textbf{(conteo)}\; $\langle g \rangle = \{ 1_{G}, g, \dots, g^{n-1} \}$ y $\lvert \langle g \rangle \rvert = \lvert g \rvert$; además $g^{i} = g^{j} \iff n \mid i - j$. \hfill\textit{Proposición 1.6.10.}
\item \textbf{(potencias)}\; $\lvert g^{k} \rvert = \dfrac{n}{\mcd(n,k)}$; en particular, si $c \mid n$, $\lvert g^{c} \rvert = n/c$. \hfill\textit{Proposición 1.7.3.}
\end{enumerate}
\end{equivbox}

\begin{equivbox}{teal}{Caracterizaciones de generadores y unidades de $\Z_m$}
Sea $G = \langle g \rangle$ con $\lvert g \rvert = n$, y sea $[a] \in \Z_{m}$:
\begin{enumerate}[label=(\arabic*),leftmargin=2em]
\item \textbf{(generador de potencia)}\; $G = \langle g^{k} \rangle \iff \mcd(n,k) = 1$. \hfill\textit{Corolario 1.7.6.}
\item \textbf{(igualdad de subgrupos)}\; $\langle g^{k} \rangle = \langle g^{\ell} \rangle \iff \mcd(k,n) = \mcd(\ell,n) \iff \lvert g^{k} \rvert = \lvert g^{\ell} \rvert$. \hfill\textit{Teorema 1.7.5.}
\item \textbf{(generador de $\Z_m$)}\; $\langle [a] \rangle = \Z_{m} \iff \MCD(a,m) = 1$. \hfill\textit{Teorema 1.6.8.}
\item \textbf{(unidad)}\; $[a] \in U(m) \iff \MCD(a,m) = 1$. \hfill\textit{Ejercicio 1.6.7.}
\end{enumerate}
\end{equivbox}

\anexogroup{Coclases y normalidad}

\begin{equivbox}{purple}{Caracterizaciones de la igualdad de coclases $aH = bH$}
Sea $H \leq G$ y $a, b \in G$. Las siguientes son equivalentes:
\begin{enumerate}[label=(\arabic*),leftmargin=2em]
\item \textbf{(criterio algebraico)}\; $b^{-1} \ast a \in H$. \hfill\textit{Lema 2.4.3.}
\item \textbf{(pertenencia)}\; $a \in bH$. \hfill\textit{Prueba del Lema 2.4.3.}
\item \textbf{(no disjuntas)}\; $aH \cap bH \neq \emptyset$. \hfill\textit{Teorema 2.4.4.}
\item \textbf{(relación de equivalencia)}\; $a \sim b$ bajo la relación cuyas clases son las coclases. \hfill\textit{Nota 2.4.5.}
\end{enumerate}
\end{equivbox}

\begin{equivbox}{orange!85!black}{Caracterizaciones de subgrupo normal $H \triangleleft G$}
Sea $H \leq G$. Las siguientes son equivalentes:
\begin{enumerate}[label=(\arabic*),leftmargin=2em]
\item \textbf{(definición)}\; $xH = Hx$ para todo $x \in G$. \hfill\textit{Definición 2.5.2.}
\item \textbf{(conjugación, igualdad)}\; $xHx^{-1} = H$ para todo $x \in G$. \hfill\textit{Nota 2.5.3.}
\item \textbf{(conjugación, contención)}\; $xHx^{-1} \subseteq H$ para todo $x \in G$. \hfill\textit{Nota 2.5.3.}
\item \textbf{(estabilidad por $C_a$)}\; $C_{a}(H) \subseteq H$ para todo $a \in G$. \hfill\textit{Ejercicio 2.6.8.}
\item \textbf{(kernel)}\; $H = \Ker(f)$ para algún homomorfismo $f$ con dominio $G$. \hfill\textit{Teoremas 2.5.5 y 2.6.4.}
\end{enumerate}
\end{equivbox}

\anexogroup{Homomorfismos y acciones}

\begin{equivbox}{brown!85!black}{Caracterizaciones de inyectividad de un homomorfismo}
Sea $f : G \to H$ un homomorfismo de grupos. Las siguientes son equivalentes:
\begin{enumerate}[label=(\arabic*),leftmargin=2em]
\item \textbf{(definición)}\; $f$ es un monomorfismo (función inyectiva). \hfill\textit{Definición 2.1.9.}
\item \textbf{(kernel trivial)}\; $\Ker(f) = \{ 1_{G} \}$. \hfill\textit{Teorema 2.1.11.}
\item \textbf{(fibras)}\; Cada fibra de $f$ es una sola coclase: $\{ x : f(x) = f(a) \} = a\Ker(f)$, que se reduce a $\{a\}$. \hfill\textit{Proposición 2.4.13.}
\end{enumerate}
\end{equivbox}

\begin{equivbox}{olive}{Caracterizaciones de una acción de grupo}
Sea $G$ un grupo y $X$ un conjunto. Dar una acción de $G$ sobre $X$ equivale a:
\begin{enumerate}[label=(\arabic*),leftmargin=2em]
\item \textbf{(axiomas)}\; Una función $\alpha : G \times X \to X$ con $\alpha(g, \alpha(h,x)) = \alpha(g \cdot h, x)$ y $\alpha(1_{G}, x) = x$. \hfill\textit{Definición 3.1.1.}
\item \textbf{(homomorfismo)}\; Un homomorfismo $\phi : G \to S_{X}$ (vía $\phi(g) = \alpha_{g}$); la correspondencia es biyectiva. \hfill\textit{Teorema 3.1.6.}
\item \textbf{(biyecciones)}\; Cada $\alpha_{g}$ es una biyección de $X$, con $(\alpha_{g})^{-1} = \alpha_{g^{-1}}$. \hfill\textit{Lema 3.1.5.}
\end{enumerate}
\end{equivbox}

\begin{equivbox}{magenta!80!black}{Paquete de Sylow: existencia, conjugación y conteo}
Sea $G$ finito, $p$ primo y $\lvert G \rvert = p^{e}k$ con $p \nmid k$:
\begin{enumerate}[label=(\arabic*),leftmargin=2em]
\item \textbf{(definición)}\; Un $p$-subgrupo de Sylow es un $H \leq G$ con $\lvert H \rvert = p^{e}$. \hfill\textit{Definición 3.4.1.}
\item \textbf{(existencia)}\; Siempre existe un $p$-subgrupo de Sylow. \hfill\textit{Teorema 3.4.3.}
\item \textbf{(conjugación)}\; Dos $p$-subgrupos de Sylow cualesquiera son conjugados. \hfill\textit{Teorema 3.4.4.}
\item \textbf{(conteo)}\; El número $n_{p}$ de $p$-subgrupos de Sylow cumple $p \mid n_{p} - 1$. \hfill\textit{Teorema 3.4.5.}
\end{enumerate}
\end{equivbox}

% =============================================================================
\section{Anexo: implicaciones y equivalencias}
% =============================================================================

Este anexo organiza, por familias de propiedades, qué implica qué: las equivalencias ($\Leftrightarrow$), las implicaciones válidas en general ($\Rightarrow$) y las recíprocas que \emph{fallan} ($\not\Rightarrow$), con contraejemplos tomados del propio curso.

\begin{equivbox}{blue}{Cíclicos, abelianos y subgrupos}
\textbf{Equivalencias ($\Leftrightarrow$):}
\begin{itemize}[leftmargin=1.8em]
\item $G$ abeliano $\iff Z(G) = G$. \hfill\textit{Nota 1.9.2.}
\item $G$ cíclico con $\lvert G \rvert = n$ $\iff G \cong \Z_{n}$; \quad cíclico infinito $\iff G \cong \Z$. \hfill\textit{Teorema 2.1.18.}
\item $G$ abeliano y simple $\iff G \cong \Z_{p}$, $p$ primo. \hfill\textit{Teorema 3.3.2.}
\end{itemize}
\textbf{Implicaciones siempre válidas ($\Rightarrow$):}
\begin{itemize}[leftmargin=1.8em]
\item $G$ cíclico $\implies$ $G$ abeliano. \hfill\textit{usado en la prueba del Teorema 3.2.14.}
\item $G$ abeliano $\implies$ todo $H \leq G$ es normal. \hfill\textit{Nota 2.5.3.}
\item $H \leq G$ cíclico $\implies$ $H$ cíclico. \hfill\textit{Teorema 1.7.7.}
\item $G$ cíclico $\implies$ $G$ numerable. \hfill\textit{Nota 1.6.9.}
\end{itemize}
\textbf{Implicaciones que NO valen en general ($\not\Rightarrow$):}
\begin{itemize}[leftmargin=1.8em]
\item abeliano $\not\Rightarrow$ cíclico: $(\R, +)$ es abeliano pero no numerable, luego no cíclico. \hfill\textit{Nota 1.6.9.}
\item grupo $\not\Rightarrow$ abeliano: $D_{4}$ y $S_{3}$ no son conmutativos. \hfill\textit{Ejercicio 1.3.8, Notación 1.8.3.}
\item asociatividad $+$ neutro $\not\Rightarrow$ grupo: $(\mathbb{N}, +)$ es monoide sin inversos. \hfill\textit{Ejemplo 1.3.2.}
\end{itemize}
\end{equivbox}

\begin{equivbox}{green!60!black}{Normalidad y cocientes}
\textbf{Equivalencias ($\Leftrightarrow$):}
\begin{itemize}[leftmargin=1.8em]
\item $H \triangleleft G \iff xHx^{-1} \subseteq H\ \forall x \iff H$ es kernel de un homomorfismo. \hfill\textit{Nota 2.5.3; Teoremas 2.5.5 y 2.6.4.}
\item $G/H$ abeliano $\iff G' \triangleleft H$ \ (con $H \triangleleft G$ y $G'$ el subgrupo conmutador). \hfill\textit{Teorema 2.6.9.}
\end{itemize}
\textbf{Implicaciones siempre válidas ($\Rightarrow$):}
\begin{itemize}[leftmargin=1.8em]
\item $K \triangleleft G$ y $K \leq H \leq G$ $\implies$ $K \triangleleft H$. \hfill\textit{Proposición 2.5.4.}
\item $N \triangleleft G$ $\implies$ $NH = HN$ para todo $H \leq G$. \hfill\textit{Teorema 2.7.3.}
\item $H \triangleleft G$ $\implies$ $G/H$ es un grupo con $aH \cdot bH = (a \ast b)H$. \hfill\textit{Teorema 2.6.2.}
\end{itemize}
\textbf{Implicaciones que NO valen en general ($\not\Rightarrow$):}
\begin{itemize}[leftmargin=1.8em]
\item $H \leq G \not\Rightarrow H \triangleleft G$: en $S_{3}$, $H = \langle (1\;2) \rangle$ tiene coclases izquierdas y derechas distintas. \hfill\textit{Ejemplos 2.4.2.}
\item $H \triangleleft G \not\Rightarrow$ $x \ast h = h \ast x$ elemento a elemento: solo garantiza $x \ast h = h' \ast x$ para algún $h' \in H$. \hfill\textit{Nota 2.5.3.}
\end{itemize}
\end{equivbox}

\begin{equivbox}{red!75!black}{Órdenes y existencia de subgrupos}
\textbf{Equivalencias ($\Leftrightarrow$):}
\begin{itemize}[leftmargin=1.8em]
\item $g^{k} = 1_{G} \iff \lvert g \rvert \mid k$. \hfill\textit{Teorema 1.6.3.}
\end{itemize}
\textbf{Implicaciones siempre válidas ($\Rightarrow$):}
\begin{itemize}[leftmargin=1.8em]
\item $H \leq G$ finito $\implies \lvert H \rvert \mid \lvert G \rvert$, y $\lvert g \rvert \mid \lvert G \rvert$. \hfill\textit{Teorema 2.4.10, Corolario 2.4.11.}
\item $p$ primo, $p \mid \lvert G \rvert$ $\implies$ existe $g \in G$ con $\lvert g \rvert = p$. \hfill\textit{Teorema 3.2.9 (Cauchy).}
\item $G$ abeliano finito y $d \mid \lvert G \rvert$ $\implies$ existe $H \leq G$ con $\lvert H \rvert = d$. \hfill\textit{Teorema 3.2.14.}
\item $\lvert G \rvert = p^{e}$ $\implies$ para todo $k \leq e$ existe $H \leq G$ con $\lvert H \rvert = p^{k}$. \hfill\textit{Teorema 3.2.15.}
\item $\lvert G \rvert = p^{2}$ $\implies$ $G$ abeliano. \hfill\textit{Corolario 3.2.13.}
\item $G$ $p$-grupo no trivial $\implies Z(G) \neq \{ 1_{G} \}$. \hfill\textit{Ejercicio 3.2.12.}
\end{itemize}
\textbf{Implicaciones que NO valen en general ($\not\Rightarrow$):}
\begin{itemize}[leftmargin=1.8em]
\item $d \mid \lvert G \rvert \not\Rightarrow$ existe \emph{elemento} de orden $d$ (si $d$ no es primo): $8 \mid 120 = \lvert S_{5} \rvert$, pero $S_{5}$ no tiene elementos de orden $8$. \hfill\textit{Nota 3.2.10.}
\item simple $\not\Rightarrow$ abeliano: $A_{5}$ es simple y no es abeliano. \hfill\textit{Teorema 3.3.5.}
\end{itemize}
\end{equivbox}

\begin{equivbox}{violet}{Permutaciones, paridad y simplicidad}
\textbf{Equivalencias ($\Leftrightarrow$):}
\begin{itemize}[leftmargin=1.8em]
\item $\alpha$ es un $r$-ciclo $\iff \lvert \alpha \rvert = r$. \hfill\textit{Ejercicio 1.8.9.}
\item $\alpha$ impar $\iff \sgn(\alpha) = -1$. \hfill\textit{Teorema 1.8.23.}
\end{itemize}
\textbf{Implicaciones siempre válidas ($\Rightarrow$):}
\begin{itemize}[leftmargin=1.8em]
\item $\alpha, \beta$ disjuntas $\implies \alpha \circ \beta = \beta \circ \alpha$. \hfill\textit{Teorema 1.8.12.}
\item Toda $\alpha \in S_{n}$ es producto de ciclos disjuntos, y de transposiciones ($n \geq 2$). \hfill\textit{Teorema 1.8.15, Proposición 1.8.16.}
\item $\sgn(\alpha\beta) = \sgn(\alpha)\sgn(\beta)$. \hfill\textit{Teorema 1.8.21.}
\item $n \geq 5 \implies A_{n}$ es simple. \hfill\textit{Teorema 3.3.7.}
\end{itemize}
\textbf{Implicaciones que NO valen en general ($\not\Rightarrow$):}
\begin{itemize}[leftmargin=1.8em]
\item conmutar $\not\Rightarrow$ regla general: dos permutaciones cualesquiera no tienen por qué conmutar ($\alpha \circ \beta \neq \beta \circ \alpha$ en $S_{3}$). \hfill\textit{Notación 1.8.3.}
\end{itemize}
\end{equivbox}

% =============================================================================
\section{Anexo: estrategias de demostración}
% =============================================================================

Patrones de argumentación que se repiten a lo largo del curso, con sus variantes y los lugares del documento donde aparecen.

\begin{strategybox}{1.\ Para mostrar que $H$ es subgrupo: criterio $x \ast y^{-1}$}
\textbf{Patrón:} \textit{verificar $H \neq \emptyset$ y que $x, y \in H \implies x \ast y^{-1} \in H$ (Teorema 1.5.2).}
\begin{itemize}[leftmargin=1.5em]
\item Para la no-vacuidad, el testigo usual es $1_{G}$.
\item Alternativa de tres pasos: cierre, neutro, inversos (Nota 1.5.1) cuando el criterio compacto enreda.
\end{itemize}
\textbf{Apariciones:} Lema 1.5.3 (intersección), Teorema 2.1.7 (kernel e imagen), Teorema 2.1.12 (imagen inversa), Lema 3.2.4 (estabilizador).
\end{strategybox}

\begin{strategybox}{2.\ Para mostrar normalidad: basta $gHg^{-1} \subseteq H$}
\textbf{Patrón:} \textit{la contención para todo $g$ ya implica la igualdad (Nota 2.5.3), así que nunca se prueba la igualdad directamente.}
\begin{itemize}[leftmargin=1.5em]
\item Si $H = \Ker(f)$, la normalidad es automática: $f(x h x^{-1}) = 1$ (Teorema 2.5.5).
\item En cocientes: $(aK)(bK)(aK)^{-1} = (aba^{-1})K$ reduce la normalidad en $G/K$ a la normalidad en $G$ (Teorema 2.7.2).
\end{itemize}
\textbf{Apariciones:} Teorema 2.5.5, Teorema 2.6.9 (conmutador), Teorema 2.7.2 ($H/K \triangleleft G/K$), Teorema 2.7.4 ($N \cap H \triangleleft H$), Ejemplo 2.6.11 ($D_4$).
\end{strategybox}

\begin{strategybox}{3.\ Para construir isomorfismos de cocientes: fabricar $f$ y aplicar el primer teorema}
\textbf{Patrón:} \textit{en lugar de construir el isomorfismo a mano, defina un homomorfismo $f$ desde el grupo grande, calcule $\Ker(f)$ e $\Imm(f)$, y aplique $G/\Ker(f) \cong \Imm(f)$ (Teorema 2.7.1).}
\begin{itemize}[leftmargin=1.5em]
\item La elección de $f$ está dictada por el cociente buscado: proyecciones $aK \mapsto aH$ (tercer teorema), inclusiones $h \mapsto hN$ (segundo teorema), $\sigma \mapsto \sgn(\sigma)$ ($S_3/A_3$).
\item Subproducto gratuito: $\Ker(f)$ resulta normal sin trabajo extra.
\end{itemize}
\textbf{Apariciones:} Teorema 2.7.2, Teorema 2.7.4, Teorema 3.1.7 (Cayley), Ejemplo 2.6.10 ($S_{3}/A_{3} \cong \Z_{2}$).
\end{strategybox}

\begin{strategybox}{4.\ Bien-definición de funciones sobre cocientes}
\textbf{Patrón:} \textit{toda función definida eligiendo representantes ($aH \mapsto \dots$) exige verificar que representantes de la misma coclase dan la misma imagen; la herramienta es el criterio $aH = bH \iff b^{-1} \ast a \in H$ (Lema 2.4.3).}
\begin{itemize}[leftmargin=1.5em]
\item Para operaciones (no solo funciones): verificar también compatibilidad en ambos factores (Nota 2.6.3).
\item El mismo chequeo aparece camuflado al probar inyectividad: $f(aH) = f(bH) \implies aH = bH$ es la dirección recíproca.
\end{itemize}
\textbf{Apariciones:} Teorema 2.6.2 (operación de $G/H$), Teorema 2.7.1 ($f'$), Teorema 2.7.2 ($f$), Teorema 2.7.6 ($\Phi$ y $f$), Teorema 2.4.7 ($f(aH) = Ha^{-1}$).
\end{strategybox}

\begin{strategybox}{5.\ Contar con Lagrange y órbita-estabilizador}
\textbf{Patrón:} \textit{traducir preguntas de existencia o tamaño a divisibilidad: $\lvert G \rvert = [G:H]\,\lvert H \rvert$ (Teorema 2.4.10) y $\lvert \mathcal{O}_{x} \rvert = [G:G_{x}]$ (Teorema 3.2.6).}
\begin{itemize}[leftmargin=1.5em]
\item Órdenes de elementos: $\lvert g \rvert \mid \lvert G \rvert$ (Corolario 2.4.11).
\item Índices encadenados: $[G:K] = [G:H][H:K]$ (Corolario 2.4.12).
\item Clases de conjugación: $\lvert x^{G} \rvert = [G : C_{G}(x)]$ (Corolario 3.2.8).
\end{itemize}
\textbf{Apariciones:} Teorema 1.7.9 (Lagrange cíclico), Corolarios 2.4.11–2.4.12, Corolario 3.2.7, Teorema 3.4.3 (cálculo final $\lvert G_{H} \rvert = p^{e}$).
\end{strategybox}

\begin{strategybox}{6.\ Inducción fuerte sobre $\lvert G \rvert$ (o sobre el exponente)}
\textbf{Patrón:} \textit{suponer el resultado para todos los grupos de orden menor y pasar a un subgrupo propio ($C_{G}(x)$, $Z(G)$, $\langle g \rangle$) o a un cociente ($G/H$, $G/Z(G)$) donde la hipótesis aplique.}
\begin{itemize}[leftmargin=1.5em]
\item El paso al cociente exige rescatar la conclusión vía el teorema de correspondencia (Teorema 2.7.6): un $S^{\ast} \leq G/H$ se levanta a $H \leq S \leq G$ con $\lvert S \rvert = \lvert S^{\ast} \rvert \cdot \lvert H \rvert$.
\item Variante combinatoria: inducción sobre el número de elementos movidos (permutaciones) o sobre la longitud $r$ de un ciclo.
\end{itemize}
\textbf{Apariciones:} Teorema 3.2.9 (Cauchy), Teorema 3.2.14, Teorema 3.2.15, Teorema 1.8.15, Proposición 1.8.16.
\end{strategybox}

\begin{strategybox}{7.\ Ecuación de clases}
\textbf{Patrón:} \textit{partir $G$ en $Z(G)$ y las clases de conjugación no triviales: $\lvert G \rvert = \lvert Z(G) \rvert + \sum_{i} [G : C_{G}(x_{i})]$, y extraer divisibilidades sobre $\lvert Z(G) \rvert$.}
\begin{itemize}[leftmargin=1.5em]
\item Si $p \mid \lvert G \rvert$ y $p$ divide cada $[G:C_{G}(x_{i})]$, entonces $p \mid \lvert Z(G) \rvert$.
\item De ahí: los $p$-grupos tienen centro no trivial (Ejercicio 3.2.12), insumo de $\lvert G \rvert = p^{2} \implies$ abeliano (Corolario 3.2.13).
\end{itemize}
\textbf{Apariciones:} Teorema 3.2.9 (Cauchy, prueba general), Ejercicio 3.2.12, Corolario 3.2.13.
\end{strategybox}

\begin{strategybox}{8.\ Inventar la acción correcta}
\textbf{Patrón:} \textit{para obtener información estructural, hacer actuar a $G$ (o a un subgrupo $P$) sobre un conjunto fabricado a medida y leer órbitas y estabilizadores.}
\begin{itemize}[leftmargin=1.5em]
\item Sobre sí mismo por traslación: Cayley, $G \hookrightarrow S_{G}$ (Teorema 3.1.7).
\item Sobre sí mismo por conjugación: clases $x^{G}$ y centralizadores (Ejemplos 3.2.5).
\item Sobre coclases $G/H$: estabilizadores conjugados $G_{xH} = xHx^{-1}$ (Ejemplos 3.2.5).
\item Sobre subconjuntos $\Omega = \{ T \subseteq G : \lvert T \rvert = p^{e} \}$ por traslación: primer teorema de Sylow (Teorema 3.4.3); las sugerencias de los teoremas 3.4.4 y 3.4.5 siguen el mismo molde.
\end{itemize}
\textbf{Apariciones:} Teorema 3.1.7, Ejemplos 3.2.5, Teoremas 3.4.3–3.4.5.
\end{strategybox}

\begin{strategybox}{9.\ (Permutaciones) Descomponer primero en ciclos disjuntos}
\textbf{Patrón:} \textit{ante cualquier afirmación sobre $\sigma \in S_{n}$, escribir su factorización en ciclos disjuntos (Teorema 1.8.15) y razonar tipo por tipo sobre la estructura cíclica.}
\begin{itemize}[leftmargin=1.5em]
\item La estructura cíclica clasifica los casos posibles (análisis exhaustivo de tipos en $A_{5}$ y $A_{6}$).
\item El signo se lee de la factorización: $\sgn(\sigma) = (-1)^{n-k}$ con $k$ ciclos (Definición 1.8.19).
\item Ciclos disjuntos conmutan (Teorema 1.8.12), así que el orden de los factores es irrelevante.
\end{itemize}
\textbf{Apariciones:} Definición 1.8.19, Ejemplo 1.8.18, Teoremas 3.3.5 y 3.3.6 (análisis por tipos).
\end{strategybox}

\begin{strategybox}{10.\ (Permutaciones) Conjugar para fabricar elementos de un subgrupo normal}
\textbf{Patrón:} \textit{si $H \triangleleft A_{n}$ y $\sigma \in H$, entonces $\alpha\sigma\alpha^{-1} \in H$ para todo $\alpha \in A_{n}$; combinando, $\alpha\sigma\alpha^{-1}\sigma^{-1} \in H$. Eligiendo $\alpha$ con astucia se produce dentro de $H$ un elemento más simple (idealmente un $3$-ciclo).}
\begin{itemize}[leftmargin=1.5em]
\item La conjugación preserva la estructura cíclica: $\alpha\sigma\alpha^{-1}$ tiene el mismo tipo que $\sigma$.
\item Una vez obtenido un $3$-ciclo en $H$, la conjugación de $3$-ciclos y la generación por $3$-ciclos (Ejercicios 3.3.4) fuerzan $H = A_{n}$.
\item Para reducir el soporte: el conmutador $\alpha\sigma\alpha^{-1}\sigma^{-1}$ mueve a lo sumo $6$ puntos si $\alpha$ es un $3$-ciclo.
\end{itemize}
\textbf{Apariciones:} Teorema 3.3.5 (tipos 5 y 6), Teorema 3.3.6 (tipo dos $3$-ciclos), Teorema 3.3.7 (reducción a $A_{6}$).
\end{strategybox}

% =============================================================================
\section{Anexo: órdenes y divisibilidad --- la herramienta de conteo}
% =============================================================================

El recurso transversal del curso: casi toda prueba de existencia o clasificación se reduce a comparar órdenes. Aquí están todas las fórmulas juntas.

\anexogroup{El orden de un elemento}

\begin{equivbox}{blue}{Reglas básicas del orden}
Sea $g \in G$ con $\lvert g \rvert = n < \infty$:
\begin{enumerate}[label=(\arabic*),leftmargin=2em]
\item $g^{k} = 1_{G} \iff n \mid k$. \hfill\textit{Teorema 1.6.3.}
\item $g^{i} = g^{j} \iff n \mid i - j$; \quad si $\lvert g \rvert = \infty$, $g^{i} = g^{j} \iff i = j$. \hfill\textit{Proposición 1.6.10.}
\item $\lvert g^{k} \rvert = \dfrac{n}{\mcd(n,k)}$; \quad si $c \mid n$, $\lvert g^{c} \rvert = n/c$; \quad si $\lvert g \rvert = rs$, $\lvert g^{r} \rvert = s$. \hfill\textit{Proposición 1.7.3, Ejercicio 1.7.2.}
\item $\lvert f(g) \rvert \mid \lvert g \rvert$ para todo homomorfismo $f$. \hfill\textit{Proposición 2.1.3.}
\item En el cociente: $(sH)^{m} = s^{m}H$, así que $\lvert sH \rvert \mid \lvert s \rvert$. \hfill\textit{prueba del Teorema 3.2.9.}
\end{enumerate}
\end{equivbox}

\anexogroup{Conteo global: Lagrange y sus corolarios}

\begin{equivbox}{green!60!black}{Lagrange e índices}
Sea $G$ finito, $K \leq H \leq G$ y $a \in G$:
\begin{enumerate}[label=(\arabic*),leftmargin=2em]
\item $\lvert G \rvert = [G:H] \cdot \lvert H \rvert$; en particular $\lvert H \rvert \mid \lvert G \rvert$. \hfill\textit{Teorema 2.4.10.}
\item $\lvert g \rvert \mid \lvert G \rvert$ para todo $g \in G$. \hfill\textit{Corolario 2.4.11.}
\item $[G:K] = [G:H] \cdot [H:K]$. \hfill\textit{Corolario 2.4.12.}
\item $\lvert aH \rvert = \lvert H \rvert$ (como cardinalidad) y el número de coclases izquierdas es igual al de derechas. \hfill\textit{Proposición 2.4.8, Teorema 2.4.7.}
\item En cíclicos: $\lvert H \rvert \mid n$ y para cada $k \mid n$ hay un único $H \leq G$ con $\lvert H \rvert = k$, a saber $\langle g^{n/k} \rangle$. \hfill\textit{Teoremas 1.7.9 y 1.7.10.}
\end{enumerate}
\end{equivbox}

\anexogroup{Conteo por acciones}

\begin{equivbox}{orange!85!black}{Órbitas, estabilizadores y clases}
Sea $G$ finito actuando sobre $X$, $x \in X$:
\begin{enumerate}[label=(\arabic*),leftmargin=2em]
\item $\lvert \mathcal{O}_{x} \rvert = [G : G_{x}]$; en particular $\lvert \mathcal{O}_{x} \rvert \mid \lvert G \rvert$. \hfill\textit{Teorema 3.2.6, Corolario 3.2.7.}
\item Las órbitas parten $X$: $X = \dot{\bigcup}\, \mathcal{O}_{x}$. \hfill\textit{Definición 3.2.1.}
\item Conjugación: $\lvert x^{G} \rvert = [G : C_{G}(x)]$ y $\lvert G \rvert = \lvert Z(G) \rvert + \sum_{i} \lvert x_{i}^{G} \rvert$. \hfill\textit{Corolario 3.2.8; prueba del Teorema 3.2.9.}
\item $\lvert S_{n} \rvert = n!$ \ y \ $\lvert A_{n} \rvert = n!/2$. \hfill\textit{Teorema 1.8.1, Definición 1.8.24.}
\end{enumerate}
\end{equivbox}

\anexogroup{Existencia de elementos y subgrupos}

\begin{equivbox}{purple}{Teoremas de existencia}
Sea $G$ finito y $p$ primo:
\begin{enumerate}[label=(\arabic*),leftmargin=2em]
\item \textbf{(Cauchy)}\; $p \mid \lvert G \rvert \implies$ existe $g$ con $\lvert g \rvert = p$. \hfill\textit{Teorema 3.2.9.}
\item \textbf{(abeliano)}\; $G$ abeliano y $d \mid \lvert G \rvert \implies$ existe $H \leq G$ con $\lvert H \rvert = d$. \hfill\textit{Teorema 3.2.14.}
\item \textbf{($p$-grupos)}\; $\lvert G \rvert = p^{e} \implies$ existe $H \leq G$ con $\lvert H \rvert = p^{k}$ para todo $k \leq e$. \hfill\textit{Teorema 3.2.15.}
\item \textbf{(Sylow)}\; $\lvert G \rvert = p^{e}k$, $p \nmid k$ $\implies$ existe $H \leq G$ con $\lvert H \rvert = p^{e}$. \hfill\textit{Teorema 3.4.3.}
\end{enumerate}
\end{equivbox}

\anexogroup{Tarjeta de bolsillo}

\begin{equivbox}{gray}{Todas las fórmulas de conteo en una vista}
\begin{itemize}[leftmargin=1.6em]
\item $g^{k} = 1_{G} \iff \lvert g \rvert \mid k$ \quad\textbullet\quad $\lvert g^{k} \rvert = \lvert g \rvert / \mcd(\lvert g \rvert, k)$
\item $\lvert \langle g \rangle \rvert = \lvert g \rvert$ \quad\textbullet\quad $G = \langle g^{k} \rangle \iff \mcd(\lvert g \rvert, k) = 1$ \quad\textbullet\quad $\langle [a] \rangle = \Z_{m} \iff \MCD(a,m) = 1$
\item $\lvert G \rvert = [G:H]\,\lvert H \rvert$ \quad\textbullet\quad $[G:K] = [G:H][H:K]$ \quad\textbullet\quad $\lvert aH \rvert = \lvert H \rvert$
\item $\lvert g \rvert \mid \lvert G \rvert$ \quad\textbullet\quad $\lvert f(g) \rvert \mid \lvert g \rvert$ \quad\textbullet\quad $\lvert sH \rvert \mid \lvert s \rvert$ en $G/H$
\item $\lvert \mathcal{O}_{x} \rvert = [G:G_{x}]$ \quad\textbullet\quad $\lvert x^{G} \rvert = [G:C_{G}(x)]$ \quad\textbullet\quad $\lvert G \rvert = \lvert Z(G) \rvert + \sum [G:C_{G}(x_{i})]$
\item $\lvert S_{n} \rvert = n!$ \quad\textbullet\quad $\lvert A_{n} \rvert = n!/2$ \quad\textbullet\quad $\lvert G/H \rvert = \lvert G \rvert / \lvert H \rvert$
\item Existencia: Cauchy ($p$ primo) \textbullet\ abeliano ($d$ arbitrario) \textbullet\ $p$-grupos ($p^{k}$, $k \leq e$) \textbullet\ Sylow ($p^{e}$ máximo)
\end{itemize}
\end{equivbox}


\end{document}
